\documentclass[11pt]{article}

% Font and language packages
\usepackage[T1]{fontenc}
\usepackage{array}
\usepackage{booktabs}
\usepackage[utf8]{inputenc}
\usepackage[english]{babel}
\usepackage[margin=.89in]{geometry} % sets all margins to 1 inch
\usepackage{verbatim} % provides the comment environment (portable replacement for comment.sty)
\usepackage{hyperref}
\usepackage{tabularx}
\usepackage[numbers]{natbib}



% Math packages
\usepackage{amsmath, amssymb, mathtools, amsthm}
\usepackage{hyperref}
\usepackage{graphicx}
\usepackage{color}

% Graphics and figures
\usepackage{epsfig}
\usepackage{tikz}
\usepackage{subcaption}
\usepackage{float}

% Additional tools
\usepackage{multicol}
\usepackage[normalem]{ulem} % normalem: keep \emph italic (do not convert emphasis to underline)
\usepackage{url}

% Theorem styles and new theorems
\theoremstyle{plain}
\newtheorem{thm}{Theorem}[section]
\newtheorem{cor}[thm]{Corollary}
\newtheorem{lem}[thm]{Lemma}
\newtheorem{prop}[thm]{Proposition}
\newtheorem{defn}[thm]{Definition}
\newtheorem{rem}[thm]{Remark}
\newtheorem{ex}[thm]{Example}
\newtheorem{aplem}[thm]{Approximation Lemma}
\newtheorem{corollary}{Corollary}[section]


\theoremstyle{definition}
\theoremstyle{remark}
\newtheorem{remark}{Remark}[section]

% Custom commands
\DeclarePairedDelimiter{\nint}\lfloor\rceil
\DeclarePairedDelimiter{\abs}\lvert\rvert
\renewcommand{\theparagraph}{\thesubsubsection.\Alph{paragraph}}
\renewcommand{\thesubparagraph}{\thesubsubsection.\Alph{paragraph}.\Roman{subparagraph}}
\newcommand{\vc}[1]{\boldsymbol{#1}}
\newcommand{\di}{\mathrm{d}}
\newcommand{\de}{\partial}
\newcommand{\eps}{\varepsilon} % redefine epsilon to varepsilon
\newcommand{\R}{\mathbb{R}}
\newcommand{\lzz}[1]{{\color{red} #1}} % consider a more descriptive name
\newcommand{\safeincludegraphics}[2][]{%
  \IfFileExists{#2}{\includegraphics[#1]{#2}}{%
    \fbox{\parbox[c][1.2in][c]{0.72\linewidth}{\centering Missing legacy figure:\\ \texttt{#2}}}%
  }%
}


\title{Inferring the Direction of Introgression from Allelic Rarefaction Statistics with Deep Learning}

\author{
    Andres del Castillo \\
    {\small Department of Mathematics} \\
    {\small Pennsylvania State University} \\
    {\small University Park, PA 16802, USA} \\ 
    \and
    Yitzchak Shmalo \\
    {\small Institute of Mathematics} \\
    {\small Hebrew University of Jerusalem} \\
    {\small Jerusalem 91904, Israel} \\
}


\date{} % Remove this line if you want to display the current date

\hypersetup{hidelinks,
  pdftitle={Inferring the Direction of Introgression from Allelic Rarefaction Statistics with Deep Learning},
  pdfauthor={Andres del Castillo and Yitzchak Shmalo}}
\begin{document}

\maketitle






\begin{abstract}
Introgression---the transfer of genetic material between diverged populations---shapes diversity and adaptation, yet inferring its \emph{direction} remains hard: classical statistics such as Patterson's $D$ and the $f$-statistics need an outgroup, lose power for subtle gene flow, and cannot say which population donated to which. We introduce DNNaic, a two-stage model that infers the magnitude and then the direction of gene flow from nine allelic-rarefaction statistics computed with PADZE, the open-source tool we release, using no outgroup. Because each simulated replicate yields many correlated feature rows, we evaluate only under a replicate-aware split, which we prove is the correct one.

On coalescent simulations, three-way direction accuracy rises from the chance level at the smallest migration rates to $99.7\%$ across an appreciable band (rate-averaged $75.6\%$), and a detect-then-orient gate separates appreciable migration from the no-migration control at ROC-AUC $0.99$. The signal lies in the features, not the architecture: a linear classifier is Bayes-optimal for direction and no deeper, kernel, or random-matrix method improves on it, whereas the more expressive network earns its keep only on the nonlinear magnitude regression, which is recoverable in rank but not in absolute scale below a Fisher-information threshold. On real data the feature pipeline and Patterson's $D$ transfer, but a frozen simulation-trained direction head is out of distribution---in \textit{Heliconius} it assigns the same class to four introgressing trios and an allopatric control alike---so a natural-panel score is not external validation without retraining on the target system. A supporting theory establishes that the rarefaction statistics are linear functionals of the joint site-frequency spectrum, that they orient gene flow where symmetric outgroup-free summaries provably cannot, and that a linear classifier is first-order optimal, with a many-direction extension to all ordered edges among more populations.
\end{abstract}

\tableofcontents

\section{Introduction}

\subsection{Detecting and orienting introgression}
Introgression—the exchange of genetic material between diverged populations—is a major force shaping diversity and adaptation. Yet, determining not only whether introgression occurred but also its magnitude (how much ancestry moved) and direction (who donated to whom) remains challenging with standard tools. Classical statistics such as Patterson’s $D$ (ABBA–BABA) have been central to detecting gene flow, using allele-sharing asymmetries among three populations and an outgroup \citep{patterson2012ancient,durand2011testing}, famously helping establish Neanderthal admixture into modern humans \citep{green2010neandertal}. However, $D$ and related $f$-statistics \citep{peter2016admixture} typically require an outgroup, offer no direct estimate of the direction of gene flow, and do not return an admixture proportion. Moreover, their power decays for ancient or subtle events and under complex demographies. Phylogenetic network approaches can model reticulation, but are often computationally intensive and still face intrinsic difficulty in resolving direction except in limited demographic regimes \citep{thawornwattana2023direction,pease2015detection,hibbins2022phylogenomic}.

Deep learning methods have increased power and spatial resolution, either by learning patterns directly from genotype matrices \citep{flagel2019unreasonable,chan2018likelihood} or by feeding population-genetic summary statistics to fully connected networks \citep{sheehan2016deep}; for reviews of supervised and deep learning in population genetics see \citet{schrider2018supervised,korfmann2023deep}. For example, \textit{genomatnn} uses convolutional neural networks (CNNs) to detect archaic adaptive introgression signals from windowed genotype images \citep{gower2021genomatnn}. \textit{IntroUNET} frames locus-wise introgression discovery as semantic segmentation, localizing introgressed tracts with a U-Net architecture \citep{ray2024introunet}. Other pipelines combine population-genetic summaries and supervised learning to highlight adaptively introgressed loci \citep{racimo2015archaic}. These approaches often require genome-wide matrices, complex architectures, and sometimes identified donor references or outgroups. In practice, they primarily detect presence or localize tracts; they generally do not provide a direct, unified readout of direction and magnitude of gene flow.

\paragraph{This work.} We introduce \textbf{DNNaic}, a deliberately simple, two-stage fully connected DNN---in the lineage of summary-statistic networks \citep{sheehan2016deep} rather than genotype-matrix CNNs---that infers the direction of introgression---and, where the signal permits, its magnitude---from a compact, interpretable feature set derived from allelic rarefaction. Our input comprises nine statistics inspired by ADZE \citep{szpiech2008adze}: allelic richness within each of three populations $(\alpha_1,\alpha_2,\alpha_3)$, private allelic richness for each $(\pi_1,\pi_2,\pi_3)$, and pairwise-private allelic richness $(\hat{\pi}_{12},\hat{\pi}_{13},\hat{\pi}_{23})$. These summarize how variation is partitioned across groups and do not require an outgroup. The pipeline first regresses the admixture proportion (migration rate $\omega_0$), then classifies the direction of gene flow using the same feature vector augmented with $\widehat{\omega}_0$.

\paragraph{Main empirical result.} Trained on coalescent simulations under a hierarchical
Wright--Fisher model, DNNaic recovers the \emph{direction} of introgression from the compact
rarefaction feature set wherever migration is appreciable. Because rarefaction produces $198$
correlated feature rows per simulated replicate, we evaluate under a strict replicate-aware protocol
in which no replicate is shared between training and test (Section~\ref{sec:leakage-aware-reanalysis});
Theorem~\ref{thm:leakage} shows that this replicate-level split is the mathematically correct one.
Under this protocol the three-way direction accuracy rises with the migration rate to $96.6\%$ at
the largest fixed rate ($2.5\times10^{-4}$) and $99.7\%$ across the appreciable band above it
($75.6\%$ averaged over all rates), falling to the chance rate as migration becomes negligible.
Operated as a detect-then-orient gate, the method identifies appreciable migration at ROC-AUC $0.99$
with a well-calibrated score (expected calibration error $0.03$) and then orients it. The results are
stable across random seeds and are reproduced by an independent estimator
(Section~\ref{subsec:true-replicate}), and they exceed a precomputed Patterson's $D$, which by
construction cannot orient gene flow. Migration-rate \emph{magnitude}, by contrast, is information-limited at small
rates---a limit we make precise with a Fisher-information bound (Theorem~\ref{thm:fisher})---so the
simulation-supported deliverable of this work is the direction of appreciable gene flow within the
stated generative families.

\paragraph{Contributions.} Our key contributions are:
\begin{enumerate}
  \item \textbf{A minimalist, interpretable feature design} based on nine allelic rarefaction statistics \citep{szpiech2008adze}, each summarized across loci by its mean, variance, and standard error---a $28$-dimensional input (the $9\times 3$ moments together with the rarefaction depth $g$) that requires no outgroup and drastically reduces input dimensionality relative to genotype-matrix CNNs.
  \item \textbf{An open-source software release, \textsc{PADZE}.} We release \textsc{PADZE}, a \texttt{pip}-installable \textsc{NumPy} package that computes allelic and private-allelic rarefaction statistics \citep{szpiech2008adze,kalinowski2004counting} directly from VCF, and use it to derive all feature statistics analyzed here. \textsc{PADZE} reproduces the original ADZE reference implementation to its $\sim$6-significant-figure print precision ($|\Delta|\approx 5\times10^{-6}$) across every statistic, depth, and population combination on the HGDP microsatellite panel of the ADZE paper, while being competitive-to-faster on that benchmark (up to $1.33\times$).
  \item \textbf{A unified two-stage DNN pipeline} that detects appreciable gene flow and then orients it within one framework, aligning the modeling steps to the structure of the inference problem.
  \item \textbf{Strong, seed-stable performance} on simulated data under replicate-aware evaluation. DNNaic resolves the direction of gene flow at \(96.6\%\) at the largest simulated rate and \(99.7\%\) across the appreciable band (\(75.6\%\) averaged over all rates), and detects appreciable migration at ROC-AUC \(0.99\) with a well-calibrated score, exceeding a precomputed Patterson's $D$, which labels direction \emph{below} the majority-class rate because it cannot orient gene flow by construction; the discriminative signal lies in the rarefaction feature set itself.
  \item \textbf{A rigorous evaluation protocol.} Because rarefaction yields $198$ correlated rows per replicate, we evaluate by ground-truth replicate identity---no replicate shared between training and test---which Theorem~\ref{thm:leakage} shows is the mathematically correct split for this data structure.
  \item \textbf{An all-edges scaling test.} We extend the feature map to four populations and classify all $4\times3=12$ ordered edges. Across $240$ independent baseline replicates, exact accuracy is $62.9\%$ (chance $8.3\%$) after training-fold-only augmentation by all $24$ population relabellings, with $76.3\%$ unordered-pair recovery and $82.5\%$ orientation conditional on the correct pair; negative controls return to chance and demographic shifts expose where frozen transfer fails.
  \item \textbf{A supporting theory package.} We prove that rarefaction statistics are linear functionals of the joint site-frequency spectrum (Theorem~\ref{thm:linear-sfs}); that introgression direction is locally identifiable from their first-order migration response (Corollary~\ref{cor:identifiability}) precisely in regimes where symmetric, outgroup-free summaries provably \emph{cannot} orient gene flow (Theorem~\ref{thm:symmetric}); that the across-locus \emph{variance} of the rarefaction statistics orients gene flow even in a degenerate regime where the first-order mean cannot (Proposition~\ref{prop:variance-orients}); that a \emph{linear} classifier on these features is first-order Bayes-optimal for direction, which explains why a logistic model matches the deep network and locates the signal in the features rather than the architecture (Proposition~\ref{prop:linear-optimal}); that $K=s(s-1)$ ordered edges incur only a $\sqrt{\log K}$ class-count cost when their minimum Fisher separation persists, with a risk-nonincreasing permutation symmetrization (Theorem~\ref{thm:many-directions}); that row-level splitting is optimistically biased (Theorem~\ref{thm:leakage}); and that magnitude estimation is information-limited at small rates via a Fisher-information bound (Theorem~\ref{thm:fisher}).
  \item \textbf{Explicit transfer limits.} The outgroup-free feature map is computationally portable across taxa, but its trained decision rule is not automatically biologically portable: demographic nuisance shifts and an allopatric \textit{Heliconius} control show that matched simulation or retraining and target-system controls are required before interpreting a natural-panel direction score.
\end{enumerate}
Taken together, DNNaic provides a fast and transparent route to inferring the direction of introgression---and, more cautiously, its magnitude---bridging the gap between traditional summary-based tests and high-dimensional deep learning detectors.


\subsection{Deep learning preliminaries}
\label{introduction_deep_neural_networks}

In classification tasks, the goal is to categorize elements of a set $S$ into one of $K$ distinct classes. Let $C(s)\in\{1,\dots,K\}$ be the correct class of $s\in S$. Given a labeled training set $T\subset S$, we aim to develop an approximate classifier to generalize accurate classification from the training set $T$ to the entire set $S$. To achieve this, we employ a Deep Neural Network (DNN) as a parameter-dependent classifier $\phi(s, \alpha)$, where $\alpha \in \mathbb{R}^\nu$ represents its parameters. The DNN produces a probability vector $(p_1(s), \dots, p_K(s))$, where $p_k(s):=\mathbb{P}[C(s)=k]$ represents the estimated probability distribution of $C(s)$.

The DNN model is structured as a composition $\rho \circ X(\cdot,\alpha)$, with $\rho$ as the softmax function (defined in \eqref{soft_max}) and $X(\cdot,\alpha)$ specified as follows. Let 
\begin{equation}
X(\cdot,\alpha)=\lambda\circ M_L(\cdot,\alpha)\circ\cdots\circ\lambda\circ M_1(\cdot,\alpha),
\end{equation}
where:
\begin{itemize}
    \item $M_k(\cdot,\alpha)$ is an affine function from $\mathbb{R}^{N_{k-1}}$ to $\mathbb{R}^{N_k}$ parameterized by a weight matrix $W_k \in \mathbb{R}^{N_k \times N_{k-1}}$ and a bias vector $\beta_k \in \mathbb{R}^{N_k}$, so that $M_k(x) = W_k \cdot x + \beta_k$.
    \item $\lambda: \mathbb{R}^m \rightarrow \mathbb{R}^m$ is a nonlinear \emph{activation function} applied componentwise. Throughout this paper $\lambda$ is the rectified linear unit (ReLU), $\lambda: x \in \mathbb{R}^m \mapsto (\max\{0,x_1\}, \dots, \max\{0,x_m\}) \in \mathbb{R}^m$, matching the network architectures used in Sections~\ref{sec:regression_model} and~\ref{sec:classification_dnn}.
    \item The softmax function $\rho$ normalizes the outputs of $X(\cdot,\alpha)$ into a probability vector over $\{1, \dots, K\}$. The components of $\rho$ are given by:
    \begin{equation}
    \label{soft_max}
    \rho(v)_i = \frac{\exp(v_i)}{\sum_{i=1}^{K} \exp(v_i)} \quad \text{for } v \in \mathbb{R}^K.
    \end{equation}
\end{itemize}

Given this neural network architecture, we train our classifier $\phi(\cdot, \alpha)$ by minimizing the cross-entropy loss, which quantifies the accuracy of our classifier over the training set and is defined as

\begin{equation}
\label{CE}
L(\alpha) = -\frac{1}{|T|} \sum_{s \in T} \log\left(\phi_{C(s)}(s, \alpha)\right).
\end{equation}


\subsubsection{Mean Squared Error Loss for Regression}
\label{mse_loss_regression}

In addition to classification tasks, DNNs are widely used for regression tasks, where the goal is to predict a continuous output. For such tasks, we often employ the Mean Squared Error (MSE) as the loss function to measure the prediction accuracy of the DNN. Let $T$ be a set of labeled data, with each element $s \in T$ associated with a true value $y(s) \in \mathbb{R}$. The MSE loss function is defined as

\begin{equation}
\label{MSE}
L_{\text{MSE}}(\alpha) = \frac{1}{|T|} \sum_{s \in T} \left( y(s) - \phi(s, \alpha) \right)^2,
\end{equation}
where $\phi(s, \alpha)$ is the predicted output of the DNN. By minimizing this loss, we aim to achieve accurate predictions for continuous targets.

This paper includes both classification and regression tasks. For classification we use fully connected networks trained with cross-entropy loss; for regression we use the same architecture but train with the mean squared error augmented by a relative-error penalty, so that small migration rates are not ignored (Section~\ref{sec:regression_model}). The two halves call for different tools, and we are explicit about it. Direction (classification) is a first-order-linear problem: a linear head is Bayes-optimal for it (Proposition~\ref{prop:linear-optimal}), so the network gives no accuracy gain and a linear model is at least as good. Magnitude (regression) is genuinely nonlinear, and there the network substantially outperforms a linear model. Table~\ref{tab:lin-vs-net} states both sides plainly. For magnitude we favour a network with an explicit linear path, which retains the linear fit and adds nonlinearity where the rate-to-feature map warrants it; for direction we use the linear head itself, because at this information limit even a linear-skip network trained end to end underperforms the pure linear model rather than falling back to it (Section~\ref{subsec:deep-kernel}). For any application to a new system we retrain or adapt on that system's own data rather than transferring a frozen model (Section~\ref{subsec:real-data}).

%Each type of task benefits from a tailored loss function, ensuring the model is optimized for the specific requirements of either discrete or continuous predictions.


\section{Methods}
\subsection{Simulations}

We conducted backwards-in-time coalescent simulations using \texttt{msprime} \citep{kelleher2016efficient} to model population histories across four populations. An ancestral population (P1234) diverged into P123 and P4 (the outgroup) $2{,}000$ generations ago; P123 then split into P12 and P3 $1{,}000$ generations ago; and P12 split into P1 and P2 $500$ generations ago. DNNaic itself uses only the three experimental populations P1, P2, and P3 and is outgroup-free; the fourth population P4 is retained solely to supply the outgroup required by the classical Patterson's-$D$ baseline against which we compare.

These coalescent simulations follow a hierarchical population-split (Wright--Fisher) demography, with an effective population size of $10{,}000$ for all populations.

\begin{figure}[H]
        \centering
        \fbox{\safeincludegraphics[width=0.65\linewidth]{figures/fig_demography.pdf}}
        \caption{The simulated four-population demography. An ancestral population splits into the three experimental populations P1, P2, P3 and an outgroup P4, at $500$, $1{,}000$, and $2{,}000$ generations before the present.}
        \label{fig:demography-generated}
\end{figure}

Continuous migration was introduced to explore three introgression scenarios among the experimental populations (P1, P2, P3), together with a no-migration control: Case A, migration from P1 into P2; Case B, from P2 into P3; Case C, from P3 into P2; and Case D, no introgression. Cases B and C are the same population pair with the flow reversed, so telling them apart is the donor-versus-recipient orientation that a symmetric statistic such as Patterson's $D$ cannot make.

\begin{figure}[H]
    \centering
    \fbox{\safeincludegraphics[width=0.8\linewidth]{figures/fig_introgression_scenarios.pdf}}
    \caption{The three directional introgression scenarios---A ($P1\!\to\!P2$), B ($P2\!\to\!P3$), C ($P3\!\to\!P2$)---and the no-migration control D. Cases B and C share a population pair with the flow reversed.}
    \label{fig:introgression-scenarios-generated}
\end{figure}

We sampled migration rates in two complementary ways. First, for each non-control direction (A, B, C) we used four fixed rates---$5\times10^{-7}$, $2.5\times10^{-6}$, $5\times10^{-5}$, and $2.5\times10^{-4}$---and generated $100$ replicates at each, giving $400$ fixed-rate replicates per direction ($1{,}200$ in total).

Second, to fill in the spectrum between the fixed rates, we drew $100$ additional rates from an exponential distribution with mean $2.5\times10^{-4}$, shared across the three non-control directions, and simulated five replicates per direction at each rate. This yields $1{,}500$ continuously-sampled replicates ($500$ per direction) spanning a biologically plausible range of low per-generation rates and lets the network learn behaviour between the fixed grid points.

\begin{figure}[H]
    \centering
    \fbox{\safeincludegraphics[width=0.65\linewidth]{figures/fig_migration_rate_distribution.pdf}}
    \caption{Distribution of simulated per-generation migration rates: four fixed rates ($5\times10^{-7}$ to $2.5\times10^{-4}$) supplemented by exponentially-distributed draws with mean $2.5\times10^{-4}$.}
    \label{fig:migration-rate-distribution-generated}
\end{figure}

For the no-migration control (direction D) we simulated $500$ replicates at zero rate. The full dataset therefore comprises $3{,}200$ replicates---$900$ each for directions A, B, and C and $500$ controls---and is archived, together with the evaluation code, at a public Zenodo record~\citep{delcastillo2026data}.

Each simulation simulated individuals with 1 Mbp with a uniform recombination rate of $1.78 \times 10^{-8}$ per site per generation. Ancestral histories were generated using sim\_ancestry, followed by the introduction of neutral mutations at a rate of $2 \times 10^{-8}$ per site. At the end of each simulation, we sampled 200 gene copies (haploid, ploidy 1) from each population, saved in .trees format and converted to .stru format for downstream PADZE analysis.


\subsection{Allelic Rarefaction Analysis with PADZE}

To generate the input features used in our downstream models, we computed allelic diversity and private allelic richness by rarefaction \citep{kalinowski2004counting}, following the private-allele estimators of \citet{szpiech2008adze} and accounting for differences in sample size between populations. This computation was carried out with \texttt{PADZE}, our NumPy implementation of these allelic-rarefaction estimators; \texttt{PADZE} is pip-installable and ships a command-line tool, so feature extraction is fully scriptable and reproducible. Our dataset consists of three populations—P1, P2, and P3—and we computed statistics at each genetic locus across the genome.

Let a locus contain $I$ distinct alleles, and let $N_{ij}$ be the number of copies of allele $i$ observed in population $j$, where $j \in \{\text{P1}, \text{P2}, \text{P3}\}$. The total number of alleles sampled at that locus in population $j$ is:

\[
N_j = \sum_{i=1}^I N_{ij}.
\]

To correct for unequal sample sizes, \texttt{PADZE} applies rarefaction by considering a standardized subsample size $g$ (same across populations). The probability that allele $i$ is \textit{not} observed in a subsample of size $g$ from population $j$ is:

\begin{equation}
Q_{ijg} = \frac{\binom{N_j - N_{ij}}{g}}{\binom{N_j}{g}}.
\end{equation}

So the probability that allele $i$ \textit{is} observed in the subsample is:

\[
P_{ijg} = 1 - Q_{ijg}.
\]

\subsubsection{Allelic Richness}

The expected number of distinct alleles observed in a sample of size $g$ from population $j$—i.e., the allelic richness—is calculated as:

\begin{equation}
{\alpha}^{(j)}_g = \sum_{i=1}^I P_{ijg}.
\end{equation}

This statistic reflects the genetic diversity within each population (P1, P2, or P3) at a fixed sampling depth.

\subsubsection{Private Allelic Richness}

To quantify population-specific diversity, \texttt{PADZE} also estimates private allelic richness: the number of alleles found only in one population and not the others. For example, for population P1, this is given by:

\begin{equation}
{\pi}^{(\text{P1})}_g = \sum_{i=1}^{I} \left[ P_{i\text{P1}g} \cdot Q_{i\text{P2}g} \cdot Q_{i\text{P3}g} \right],
\end{equation}
and similarly for P2 and P3. This formula computes the expected number of alleles present in P1 but absent in both P2 and P3 when sampling $g$ alleles from each.

\subsubsection{Combination-Private Allelic Richness}

\texttt{PADZE} also estimates alleles that are exclusive to pairs of populations. For instance, alleles shared by P1 and P2 but absent from P3 are captured by:

\begin{equation}
\hat{\pi}^{(\text{P1,P2})}_{g} = \sum_{i=1}^{I} \left[ P_{i\text{P1}g} \cdot P_{i\text{P2}g} \cdot Q_{i\text{P3}g} \right].
\end{equation}

This helps identify alleles that are private not to a single population, but to specific groupings—e.g., regionally shared but geographically restricted variation. With our three-population assumption, there are exactly three possible pairwise combinations: $\hat{\pi}^{(\text{P1,P2})}_{g}$, $\hat{\pi}^{(\text{P1,P3})}_{g}$, and $\hat{\pi}^{(\text{P2,P3})}_{g}$. The other possible combinations ($\binom{3}{1}$(single populations) or $\binom{3}{3}$ (all three populations together)) either reduce to private allelic richness or allelic richness, and are thus trivial.

\subsubsection{Defining Model Inputs: the Allelic Richness Vector}

Each of the statistics described above is computed separately at every locus. Thus, to summarize the information across the genome, \texttt{PADZE} reports the following three values for each standardized sample size $g$:

(i) The mean across loci:
    \begin{equation}
    \mu_{x^{(j)}_g} = \frac{1}{L} \sum_{\ell=1}^{L} ({x^{(j)}_g})_\ell,
    \end{equation}
    where $({x^{(j)}_g})_\ell$ is the value of the statistic $x$ for population $j$ and sample size $g$ at locus $\ell$, and $L$ is the total number of loci.
    
(ii) The variance across loci:
    \begin{equation}
    \sigma^2_{x^{(j)}_g} = \frac{1}{L - 1} \sum_{\ell=1}^{L} \left(({x^{(j)}_g})_\ell - \mu_{x^{(j)}_g} \right)^2,
    \end{equation}
    
(iii) The standard error:
    \begin{equation}
    \text{SE}_{x^{(j)}_g} = \frac{\sigma_{x^{(j)}_g}}{\sqrt{L}}.
    \end{equation}


From these summary statistics, we can define the allelic richness vector for a specific standardized sample size 'g' to be used as the input for our downstream models:

\[
\mathbf{x}_g =
\left\langle
g,\,
\alpha^{*(\text{P1})}_g,\,
\alpha^{*(\text{P2})}_g,\,
\alpha^{*(\text{P3})}_g,\,
\pi^{*(\text{P1})}_g,\,
\pi^{*(\text{P2})}_g,\,
\pi^{*(\text{P3})}_g,\,
\hat{\pi}^{*(\text{P1,P2})}_g,\,
\hat{\pi}^{*(\text{P1,P3})}_g,\,
\hat{\pi}^{*(\text{P2,P3})}_g
\right\rangle
\]

\vspace{0.5em}

\noindent Each term of the form \( x^{*(j)}_g \) represents a grouping of the mean, variance, and standard error of statistic \( x \) for population or population pair \( j \) at rarefaction depth \( g \):
\[
x^{*(j)}_g := \left( \mu_{x^{(j)}_g},\ \sigma^2_{x^{(j)}_g},\ \mathrm{SE}_{x^{(j)}_g} \right)
\]

In our simulation study, we sampled 200 gene copies (haploid) at the end of each replicate from each of the three populations. Thus, running a PADZE analysis on one replicate results in 198 individual allelic richness vectors, one for each standardized sample size $g \in [2, 199]$:

\[
\begin{aligned}
\mathbf{x}_2 &= \left\langle 2,\ \alpha^{*(\text{P1})}_2,\ \alpha^{*(\text{P2})}_2,\ \alpha^{*(\text{P3})}_2,\ \pi^{*(\text{P1})}_2,\ \pi^{*(\text{P2})}_2,\ \pi^{*(\text{P3})}_2,\ \hat{\pi}^{*(\text{P1,P2})}_2,\ \hat{\pi}^{*(\text{P1,P3})}_2,\ \hat{\pi}^{*(\text{P2,P3})}_2 \right\rangle \\
\mathbf{x}_3 &= \left\langle 3,\ \alpha^{*(\text{P1})}_3,\ \alpha^{*(\text{P2})}_3,\ \alpha^{*(\text{P3})}_3,\ \pi^{*(\text{P1})}_3,\ \pi^{*(\text{P2})}_3,\ \pi^{*(\text{P3})}_3,\ \hat{\pi}^{*(\text{P1,P2})}_3,\ \hat{\pi}^{*(\text{P1,P3})}_3,\ \hat{\pi}^{*(\text{P2,P3})}_3 \right\rangle \\
\mathbf{x}_4 &= \left\langle 4,\ \alpha^{*(\text{P1})}_4,\ \alpha^{*(\text{P2})}_4,\ \alpha^{*(\text{P3})}_4,\ \pi^{*(\text{P1})}_4,\ \pi^{*(\text{P2})}_4,\ \pi^{*(\text{P3})}_4,\ \hat{\pi}^{*(\text{P1,P2})}_4,\ \hat{\pi}^{*(\text{P1,P3})}_4,\ \hat{\pi}^{*(\text{P2,P3})}_4 \right\rangle \\
&\vdots \\
\mathbf{x}_{199} &= \left\langle 199,\ \alpha^{*(\text{P1})}_{199},\ \alpha^{*(\text{P2})}_{199},\ \alpha^{*(\text{P3})}_{199},\ \pi^{*(\text{P1})}_{199},\ \pi^{*(\text{P2})}_{199},\ \pi^{*(\text{P3})}_{199},\ \hat{\pi}^{*(\text{P1,P2})}_{199},\ \hat{\pi}^{*(\text{P1,P3})}_{199},\ \hat{\pi}^{*(\text{P2,P3})}_{199} \right\rangle
\end{aligned}
\]

\subsection{The PADZE Software Package}
\label{subsec:padze-software}

The feature statistics in this study were computed with \textsc{PADZE}~\citep{padze2026}, an
open-source, \texttt{pip}-installable \textsc{NumPy} implementation of allelic and
private-allelic rarefaction that we release with this work. \textsc{PADZE} computes
standardized-depth allelic richness, private allelic richness, and generalized
combination-private richness \citep{szpiech2008adze}, all built on the rarefaction
estimator for the expected number of distinct (or privately held) alleles at a fixed
subsample depth \citep{kalinowski2004counting}. It retains the classical ADZE mean,
variance, and standard-error outputs; only these three moments (mean, variance, and standard
error) enter the feature set used here. The package reads VCF directly (as well as
the original STRUCTURE format, for parity testing), ships a documented Python API and a
command-line tool, and requires only \textsc{Python} and \textsc{NumPy}, with no
compilation step.

We validated \textsc{PADZE} against the reference implementation---the original ADZE
program of \citet{szpiech2008adze}---on the Rosenberg et al.\ (2005) HGDP microsatellite
panel used in that paper's empirical example, over all $2^{5}-1=31$ region combinations
and every rarefaction depth $g=2,\dots,48$. Across every statistic, depth, and population
combination the two programs agree to the reference's $\sim$6-significant-figure print
precision (maximum absolute difference $|\Delta|\approx 5\times10^{-6}$;
Table~\ref{tab:padze-validation}), the remaining gap being only the reference's printed
rounding; the \textsc{PADZE} outputs are furthermore byte-identical across CPU
architectures (arm64 vs.\ x86-64). On the paper's exact computation, \textsc{PADZE} is
competitive-to-faster than the compiled reference---$3.96$\,s vs.\ $5.28$\,s on the H952
panel, a decisive $1.33\times$ speedup from \textsc{NumPy} vectorization of the deep
per-combination rarefaction loop. We characterize the trade-off transparently: in a
complementary regime of many populations at shallow depth, where the $2^{J}-1$
combination count dominates but each combination carries little work, the compiled
reference remains ahead at every $J$ tested (through $J=10$), with \textsc{PADZE}
monotonically closing the gap. The deciding axis is thus the per-combination workload
(rarefaction depth $\times$ sample size), not the number of populations $J$ alone: deep
rarefaction over large samples favors \textsc{PADZE}'s vectorization, exactly the setting
of the analyses in this paper.

\begin{table}[t]
\centering
\caption{Validation of \textsc{PADZE} against the reference ADZE implementation
\citep{szpiech2008adze} on the Rosenberg et al.\ (2005) HGDP microsatellite panel, over
all $31=2^{5}-1$ region combinations and depths $g=2,\dots,48$ ($1{,}692$ compared rows
per panel; $721$ loci throughout). Agreement is reported as the maximum absolute
difference in the mean, variance, and standard error over all statistics, depths, and
combinations; timing is the median wall-clock over $N=12$ interleaved reps on an isolated
2-vCPU CPU VM. \textsc{PADZE} matches the reference to its print precision and is faster
on this workload.}
\label{tab:padze-validation}
\begin{tabular}{@{}lccc@{}}
\toprule
Panel & Max $|\Delta|$ & Reference (s) & \textsc{PADZE} (s) \\
\midrule
H952  & $5.0\times10^{-6}$ & $5.28$ & $\mathbf{3.96}$ \\
H1048 & $5.0\times10^{-6}$ & $4.55$ & $\mathbf{4.05}$ \\
\midrule
\multicolumn{2}{@{}l}{Speedup (\textsc{PADZE} vs.\ reference)} & \multicolumn{2}{c}{$1.33\times$ (H952), $1.12\times$ (H1048)} \\
\multicolumn{2}{@{}l}{Median peak memory} & $\sim\!52$\,MB & $\sim\!48$\,MB \\
\multicolumn{2}{@{}l}{Cross-architecture agreement (arm64 vs.\ x86-64)} & \multicolumn{2}{c}{byte-identical ($\Delta=0$)} \\
\bottomrule
\end{tabular}
\end{table}


\section{Main Results}
\label{sec:main-results}

\paragraph{Principal findings.} DNNaic resolves the \emph{direction} of appreciable introgression. Under strict replicate-aware evaluation (no simulation replicate shared between training and test), three-way direction accuracy rises with the migration rate to $96.6\%$ at the largest fixed rate ($2.5\times10^{-4}$) and $99.7\%$ across the appreciable band above it ($75.6\%$ averaged over all rates), and falls to the chance rate as migration becomes negligible. Operated as a detect-then-orient gate, the method identifies appreciable migration at ROC-AUC $0.99$ with a well-calibrated score (expected calibration error $0.03$), then orients it. A separate four-population experiment resolves all 12 ordered edges at $62.9\%$ exact accuracy versus $8.3\%$ chance when the known permutation symmetry is enforced, but frozen accuracy degrades under weaker signal and mismatched nuisance distributions. These results are stable across random seeds and exceed a precomputed Patterson's $D$, which by construction cannot orient gene flow. All rarefaction features are computed with \textbf{PADZE}, our open-source, pip-installable tool released with this work and validated against the reference ADZE implementation \citep{szpiech2008adze}.

\paragraph{Idea (two-stage pipeline).}
We model migration in two steps using the same allelic richness statistics vector as input: (i) a regression DNN first predicts the magnitude $\omega_0$; (ii) a classification DNN then predicts the direction $\beta_0$, taking as features the same allelic richness vector together with the predicted magnitude $\widehat{\omega}_0$. This decouples scale from sign/orientation and empirically stabilizes learning.

\subsection{Metrics and definitions}
\label{subsec:metrics}
We report standard losses and a scale-aware relative error:

\begin{itemize}
    \item \textbf{Mean Squared Error (MSE).} For scalar target $y$ and prediction $\hat y$, $\mathrm{MSE}=\frac{1}{n}\sum_{i=1}^n (y_i-\hat y_i)^2$.
    \item \textbf{Mean Relative Error (MRE, \%)} (``relative loss''): 
    \begin{equation}
      \label{eq:mre}
      \mathrm{MRE}(\%) \;=\; 100 \times \frac{1}{n}\sum_{i=1}^n 
      \frac{|y_i-\hat y_i|}{\max(\varepsilon,\;|y_i|)}\,,
    \end{equation}
    where $\varepsilon>0$ (we take $\varepsilon=10^{-8}$) prevents division by zero when some $y_i$ is (near) zero.\footnote{If targets are strictly nonzero, \eqref{eq:mre} reduces to the usual mean absolute percentage error.}
    \item \textbf{Classification accuracy / loss.} Accuracy is the fraction correct; loss is the average cross-entropy on the validation set.
\end{itemize}

\subsection{Benchmark Results}
\label{subsec:benchmarks}

We evaluate every reported result under strict replicate-aware splitting, because the
alternative is badly misleading. A naive same-replicate split---which lets rarefaction curves
from one replicate appear in both train and test---inflates apparent direction accuracy to
$\approx99.6\%$, an artifact of leakage rather than genuine generalization
(Section~\ref{sec:leakage-aware-reanalysis}, Theorem~\ref{thm:leakage}). The leakage-controlled
estimates reported in Section~\ref{sec:leakage-aware-reanalysis}, obtained by splitting on
ground-truth replicate identity, are the numbers we treat as real.

\subsection{Learning curves and qualitative behavior}
\label{subsec:learning-curves}

\begin{figure}[H]
  \centering
  \fbox{\safeincludegraphics[width=0.9\linewidth]{figures/fig_pipeline.pdf}}
  \caption{\textbf{Proposed two-stage DNN pipeline.} The allelic richness statistics vector feeds a regression head to predict the migration magnitude $\omega_0$. The predicted magnitude $\widehat{\omega}_0$, concatenated with the original allelic features, then feeds a classification head to predict the migration direction $\beta_0$. This design mirrors the problem structure (scale first, orientation second) and empirically reduces classification ambiguity.}
  \label{fig:pipeline-detailed-generated}
\end{figure}

\begin{figure}[H]
  \centering
  \begin{subfigure}[b]{0.49\textwidth}
    \centering
    \fbox{\safeincludegraphics[width=\linewidth]{figures/fig_reg_learning.pdf}}
    \caption{\textbf{Regression training loss over epochs} for $\omega_0$}
    \label{fig:reg-train-loss-detailed}
  \end{subfigure}
  \hfill
  \begin{subfigure}[b]{0.49\textwidth}
    \centering
    \fbox{\safeincludegraphics[width=\linewidth]{figures/fig_real_array_stress_tests.pdf}}
    \caption{\textbf{Magnitude regression metrics.} The regression head optimizes stably; magnitude is reliably recovered only in the appreciable regime, consistent with the information limit of Theorem~\ref{thm:fisher}.}
    \label{fig:real-array-stress-tests-generated}
  \end{subfigure}
  \caption{\textbf{Regression DNN (migration magnitude $\omega_0$).} Learning curves demonstrate fast optimization, followed by a stable generalization regime without overfitting.}
  \label{fig:regression-dnn-combined}
\end{figure}

\begin{figure}[H]
  \centering
  \begin{subfigure}[b]{0.49\textwidth}
    \centering
    \fbox{\safeincludegraphics[width=\linewidth]{figures/fig_clf_learning.pdf}}
    \caption{\textbf{Classification training loss} for direction classification of $\beta_0$}
    \label{fig:class-train-loss-detailed}
  \end{subfigure}
  \hfill
  \begin{subfigure}[b]{0.49\textwidth}
    \centering
    \fbox{\safeincludegraphics[width=\linewidth]{figures/fig_confusion_leakagefree.pdf}}
    \caption{\textbf{Leakage-controlled direction confusion} (three-way, replicate-aggregated, all positive rates). The residual confusion concentrates on the sister-to-sister scenario A---the weakest-footprint event---while the reversed-flow pair B/C is the best separated (pair confusions $307\,{:}\,217\,{:}\,160$ for A--B, A--C, B--C): the asymmetric channels orient the reversed pair precisely where a symmetric $D$ cannot (Theorem~\ref{thm:symmetric}).}
    \label{fig:class-confusion-detailed}
  \end{subfigure}
  \caption{\textbf{Classification DNN (migration direction $\beta_0$).} (a) Training loss over epochs (held-in dynamics). (b) The leakage-controlled test confusion by ground-truth replicate identity (Section~\ref{sec:leakage-aware-reanalysis}); orientation of the reversed-flow pair B/C---the donor-versus-recipient call a symmetric $D$ cannot make (Theorem~\ref{thm:symmetric})---is in fact the cleanest, and the residual errors fall on the low-footprint sister-to-sister scenario A, whose first-order signal is the weakest (Remark~\ref{rem:threshold-numerics}).}
  \label{fig:classification-dnn-combined}
\end{figure}

\subsection{What the results say}
\label{subsec:interpretation}
The direction signal in the rarefaction statistics is, to first order, \emph{linear}: a logistic
model on the features matches the network (Section~\ref{subsec:true-replicate}), exactly as
Proposition~\ref{prop:linear-optimal} predicts. The two-stage design---estimating magnitude first,
then feeding $\widehat{\omega}_0$ into the direction classifier---sharpens class separation at larger
rates but is not what creates the signal. The core capability is directional: under
leakage-controlled evaluation the classifier recovers introgression direction with high accuracy
where migration is appreciable. Magnitude regression is intrinsically harder and
degrades at small rates (Section~\ref{sec:regression_model}), an information limit made precise by
the Fisher-information bound of Theorem~\ref{thm:fisher}.

\subsection{Reporting details}
\label{subsec:reporting}
All reported numbers use leakage-controlled splits by ground-truth replicate identity
(Section~\ref{sec:leakage-aware-reanalysis}); the DNN regression and classifier use fixed
training horizons with the same allelic richness vector feeding both stages, and model
comparisons use identical splits and preprocessing. Throughout, MRE (\%) is the scale-aware
relative loss of \eqref{eq:mre} and classification ``loss'' denotes cross-entropy.

% ---- Optional bridging remark to theory (if desired) ----



\section{Feature geometry and supporting analyses}
\label{sec:supporting}

\subsection{Exploratory geometry of the rarefaction features}
\label{sec:eda}
As a first check on the discriminatory power of the rarefaction features, we ran principal component analysis (PCA) on the standardized feature vectors---summary statistics (mean, variance, and standard error) of the nine allelic richness measures computed by PADZE at a standardized sample size $g$, giving $28$ features per replicate. Under strong gene flow (an exaggerated $25\%$ migration rate used only for this visualization) the four scenarios separate hierarchically in PC space, mirroring the population tree of Figure~\ref{fig:demography-generated}: the no-migration control (Case D) and Case A ($P1\rightarrow P2$) stand apart, while the recently diverged pair $P2$/$P3$ (Cases B and C) overlaps. At the realistic rates present in our study the separation shrinks but does not vanish---population structure and the tree-like PC1--PC2 geometry persist even as lower and mixed migration rates blur the class boundaries (Figure~\ref{fig:pca-comparison}), motivating the supervised, replicate-aware treatment developed in the rest of the paper.

\begin{figure}[H]
  \centering
  \begin{subfigure}[b]{0.49\textwidth}
    \centering
    \fbox{\safeincludegraphics[width=\linewidth]{figures/fig_pca_by_class.pdf}}
    \caption{PCA of downloaded processed arrays colored by class.}
    \label{fig:pca-by-class-generated-a}
  \end{subfigure}
  \hfill
  \begin{subfigure}[b]{0.49\textwidth}
    \centering
    \fbox{\safeincludegraphics[width=\linewidth]{figures/fig_pca_by_rate.pdf}}
    \caption{PCA of downloaded processed arrays colored by migration rate.}
    \label{fig:pca-by-rate-generated-a}
  \end{subfigure}
  \caption{3D PCA analysis of allelic richness vectors for a high migration rate and a realistic migration rate found in our dataset.}
  \label{fig:pca-comparison}
\end{figure}

\subsection{Magnitude estimation is information-limited at small rates}
\label{sec:regression_model}

The first stage of the pipeline is a regression head that maps the standardized rarefaction features to the (log) migration magnitude $\omega_0$: a four-layer elastic-net-regularized ReLU network trained with AdamW under a scheduled learning rate and a custom loss that augments the mean squared error with a relative-error penalty so that small rates are not ignored. Magnitude is the hard side of this problem, and we report it plainly. Under the strict
replicate-aware split the relative error is dominated by the smallest rates and grows without bound
as the rate falls---from about $34\%$ in the appreciable band to several thousand percent at
$5\times10^{-7}$---so absolute scale is not recoverable at small rates. Two things survive. A
median-regression (absolute-error) head restricted to the appreciable regime recovers the magnitude
to a mean relative error of about $31\%$ (median $27\%$); and the \emph{rank} of the rate is
recoverable across the whole range (Spearman correlation $\approx 0.8$ between predicted and true
rate) even where the scale is not, so the regressor still orders replicates by migration intensity
below the threshold. This is an information limit rather than an optimization failure:
Theorem~\ref{thm:fisher} gives a Fisher-information lower bound on the relative error that diverges as
$\omega_0\to 0$ at fixed genome size, and Corollary~\ref{cor:appreciable-threshold} locates the
threshold below which no estimator---DNN or otherwise---can recover scale. The predicted $\widehat{\omega}_0$ is nonetheless a useful auxiliary feature for the downstream direction classifier, which is where the discriminative signal of the paper lies.

\section{Classification DNN for Migration Direction Inference}
\label{sec:classification_dnn}

In this section, we describe a deep neural network (DNN) designed to classify objects into one of four migration groups, labeled as \(A\), \(B\), \(C\), and \(D\). Let 
\[
\mathcal{D} = \{(x_i, y_i) \mid i=1,\dots, N\}
\]
denote the training dataset, where each input \(x_i \in \mathbb{R}^d\) is a vector of \emph{standardized direction features} (denoted by \texttt{std\_dir\_features}) and the corresponding target \(y_i\) is a one-hot encoded vector in \(\{0,1\}^4\) (denoted by \texttt{OHEdir\_target}). We define a label function \(C: \mathcal{D} \to \{A, B, C, D\}\) that assigns each object to one of the four migration groups.

\subsection{Input Feature Weighting}

To emphasize features that are deemed more informative (e.g., migration rate, standardized sample size 'g', and statistics related to \(\hat{\pi}\)), we introduce a weight vector 
\[
w \in \mathbb{R}^d,
\]
defined by
\[
w_j = 
\begin{cases}
5, & \text{if } j\in\{1,2\} \cup \{d-8, \dots, d\},\\[1mm]
1, & \text{otherwise}.
\end{cases}
\]
The weighted input is then given by
\[
\tilde{x}_i = w \odot x_i,
\]
where \(\odot\) denotes the elementwise (Hadamard) product.

\subsection{Network Architecture}

The classification DNN is a function
\[
f: \mathbb{R}^d \to \Delta^3,
\]
mapping the weighted input \(\tilde{x}_i\) to a probability vector in the 4-simplex \(\Delta^3 = \{p\in\mathbb{R}^4: p_k\ge0,\; \sum_{k=1}^4 p_k = 1\}\). The architecture is defined as a sequence of fully connected layers:
\[
\begin{aligned}
h^{(1)} &= \varphi\bigl(W_1\,\tilde{x}_i + b_1\bigr) \quad \in \mathbb{R}^{300},\\[1mm]
h^{(2)} &= \varphi\bigl(W_2\,h^{(1)} + b_2\bigr) \quad \in \mathbb{R}^{300},\\[1mm]
h^{(3)} &= \varphi\bigl(W_3\,h^{(2)} + b_3\bigr) \quad \in \mathbb{R}^{300},\\[1mm]
\hat{y}_i &= \rho\bigl(W_4\,h^{(3)} + b_4\bigr) \quad \in \Delta^3,
\end{aligned}
\]
where:
\begin{itemize}
    \item \(W_k\) and \(b_k\) for \(k=1,2,3\) are the weight matrices and bias vectors for the hidden layers with dimensions \(W_k \in \mathbb{R}^{300 \times N_{k-1}}\) (with \(N_0 = d\)) and \(b_k\in \mathbb{R}^{300}\).
    \item \(\varphi(z)=\max\{0,z\}\) is the ReLU activation function.
    \item \(W_4 \in \mathbb{R}^{4\times300}\) and \(b_4 \in \mathbb{R}^{4}\) define the output layer.
    \item \(\rho: \mathbb{R}^4 \to \Delta^3\) is the softmax function given by
    \[
    \rho(v)_k = \frac{\exp(v_k)}{\sum_{j=1}^{4}\exp(v_j)}, \quad k=1,\dots,4.
    \]
\end{itemize}

Additionally, elastic net regularization is applied to the weights of the hidden layers. That is, for each hidden layer we incorporate a penalty of the form
\[
\lambda_1 \|W_k\|_1 + \lambda_2 \|W_k\|_2^2,
\]
with small coefficients \(\lambda_1 = \lambda_2 = 10^{-7}\).

\subsection{Training Objective and Learning Rate Schedule}

The network is trained by minimizing the categorical cross-entropy loss over the training set:
\[
L(\alpha) = -\frac{1}{N}\sum_{i=1}^{N} \sum_{k=1}^{4} y_{ik}\log\bigl(\hat{y}_{ik}\bigr),
\]
where \(y_{ik}\) denotes the \(k\)-th component of the one-hot encoded target \(y_i\), and \(\alpha\) represents the collection of all parameters in the network.

An AdamW optimizer is used with an initial learning rate of \(10^{-3}\). A learning rate schedule is defined by a piecewise function:
\[
\eta(t)=
\begin{cases}
9\times 10^{-4}, & t < 15, \\[1mm]
-\bigl(1.274\times 10^{-5}\bigr) t + 0.001091, & 15 \le t < 85,
\end{cases}
\]
where \(t\) denotes the epoch number. This schedule is implemented via a custom callback during training.

\subsection{Model Evaluation and Checkpointing}

During training, the dataset is randomly partitioned into training and testing subsets. At each epoch, the network's performance is monitored by evaluating both training and validation accuracy. A custom callback is used to:
\begin{enumerate}
    \item Log the accuracy metrics.
    \item Sample a few test inputs and print both the predicted and true labels.
\end{enumerate}

Finally, the trained model is checkpointed for later evaluation.

\section{Direction inference under replicate-aware evaluation}
\label{sec:leakage-aware-reanalysis}

Allelic rarefaction produces $198$ correlated feature rows per simulated replicate, so a
row-level train/test split can place rows from the same genealogy in both partitions and thereby
measure within-replicate interpolation rather than generalization to new replicates.
Theorem~\ref{thm:leakage} shows that the mathematically correct estimand requires splitting on the
replicate. We therefore evaluate DNNaic by ground-truth replicate identity: the $28$-dimensional
PADZE feature vectors are grouped by their originating replicate and split with
$\mathrm{GroupKFold}(5)$, so no simulated genealogy is ever shared between training and test. The
benchmark below uses $3{,}200$ true replicate groups---direction classes A/B/C together with a
zero-migration control D---with per-replicate probability aggregation and three independent random
seeds. The feature arrays, replicate labels, and evaluation code are archived at a public Zenodo
record~\citep{delcastillo2026data}.

\subsection{Direction and detection of appreciable gene flow}
\label{subsec:true-replicate}

Under this replicate-aware protocol DNNaic orients gene flow accurately wherever migration is
strong enough to leave a trace, and degrades gracefully to chance where it is not. Three-way
direction accuracy climbs monotonically with the migration rate: it is $34\%$ at $5\times10^{-7}$
and $39\%$ at $2.5\times10^{-6}$ (both near the $33\%$ chance level), $56\%$ at $5\times10^{-5}$,
and $96.6\%$ at the largest fixed rate $2.5\times10^{-4}$ for the boosted model
(Figure~\ref{fig:direction-by-rate}, Table~\ref{tab:true-replicate}); the Fisher-optimal linear head
is stronger still at the informative rates, $71.6\%$ at $5\times10^{-5}$ and $100\%$ at that largest
fixed rate.
Averaged over the strong-migration band---rate $\ge 2.5\times10^{-4}$, which we call
\emph{appreciable} and which also collects the higher continuously-sampled rates---the accuracy is
$99.7\%$ (three-seed mean, standard deviation $0.1\%$); averaged over all positive rates it is
$75.6\%$. The steep rise near $2.5\times10^{-4}$ is not a tuning artifact but the directional
shadow of the Fisher threshold $m_\star$ (Corollary~\ref{cor:appreciable-threshold}): below it the
first-order signal $m\Delta_\beta$ that separates directions (Proposition~\ref{prop:linear-optimal})
sinks beneath the sampling noise, and the near-zero control becomes indistinguishable from the
smallest rates.

\begin{figure}[H]
  \centering
  \safeincludegraphics[width=0.6\linewidth]{figures/fig_direction_by_rate.pdf}
  \caption{\textbf{Direction accuracy rises to the Fisher threshold.} Three-way direction accuracy
  (leakage-controlled, three-seed mean) climbs monotonically with the per-generation migration rate,
  from the $33\%$ chance level at the smallest rates to $96.6\%$ at the largest fixed rate; the shaded
  band marks the appreciable regime (rate $\ge 2.5\times10^{-4}$). The decline below the threshold is
  the directional shadow of the Fisher-information limit (Corollary~\ref{cor:appreciable-threshold}).}
  \label{fig:direction-by-rate}
\end{figure}

\begin{table}[H]
\centering
\scriptsize
\begin{tabular}{lcc}
\toprule
migration rate & boosted / network (\%) & linear (\%) \\
\midrule
$5\times10^{-7}$   & 34.0 & 29.4 \\
$2.5\times10^{-6}$ & 39.3 & 41.0 \\
$5\times10^{-5}$   & 55.8 & \textbf{71.6} \\
$2.5\times10^{-4}$ (largest fixed rate) & 96.6 & \textbf{100.0} \\
\midrule
appreciable band ($\ge 2.5\times10^{-4}$) & \textbf{99.7} & \textbf{100.0} \\
all positive rates & 75.6 & \textbf{78.1} \\
\bottomrule
\end{tabular}
\caption{\textbf{Three-way direction accuracy under the replicate-aware split}
($\mathrm{GroupKFold}(5)$ on $3{,}200$ true replicate groups; three-way chance $33\%$). The
boosted/network column is the three-seed figure of Section~\ref{subsec:benchmarks}; the linear column
is a logistic model under the identical leakage-free protocol. Accuracy rises monotonically with the
migration rate, and the Fisher-optimal linear head (Proposition~\ref{prop:linear-optimal}) is
stronger at the informative rates---$71.6\%$ versus $55.8\%$ at $5\times10^{-5}$ and $100\%$ at the
largest fixed rate---while both fall to the chance level as migration becomes negligible, the
directional shadow of the Fisher threshold (Corollary~\ref{cor:appreciable-threshold}). The
appreciable band also collects the higher continuously-sampled rates.}
\label{tab:true-replicate}
\end{table}

Operated as a \emph{detect-then-orient} gate---first flag whether migration is appreciable, then
orient it only where detection succeeds---the method is accurate and well calibrated within the
evaluated simulation distribution. The
detection gate reaches ROC-AUC $0.99$ separating appreciable migration from the no-migration
control ($0.80$ across all rates), with expected calibration error $0.03$ and Brier score $0.11$,
so its confidence scores can be read directly to set an operating point; against the more
demanding mixture of weak-positive and control replicates it reaches $0.96$
(\texttt{scripts/appreciable\_gate.py}). The direction head supports the same selective use: ranking
appreciable replicates by the classifier's confidence and abstaining on the least-confident quarter
raises three-way accuracy from $99.6\%$ to $99.9\%$, and abstaining on the least-confident half makes
it perfect on the replicates retained. This is the usable form of the method: it asserts a direction
only where the signal is present and abstains where it is not.

Two checks confirm the result is a property of the data rather than of a single model. The accuracy
is stable across all three seeds (standard deviations at or below $0.3\%$), and an independent
estimator built from a different software stack---an \texttt{sklearn} gradient-boosted classifier
under the identical leakage-free protocol---reproduces the figures to within a point (appreciable
gate AUC $0.99$, appreciable direction $99.8\%$, overall $74.7\%$, and the same per-rate decay).
A third check strengthens the leakage control from cross-validation within one simulation set to a
physically separate test set: we train on the canonical $3{,}200$-replicate benchmark and evaluate
on two further batches of $3{,}200$ replicates each, generated in independent runs with freshly
drawn migration rates and disjoint seeds, so that no genealogy, rate draw, or replicate is shared
with training. On these unseen batches the estimator reproduces the within-set figures---rate-averaged
direction $76.5\%$, appreciable $99.7\%$, detection gate ROC-AUC $0.81$ overall and $0.99$ in the
appreciable band, expected calibration error $0.02$ (linear head; the nonlinear head is within a
point), and pooling the two batches into a single $6{,}400$-replicate held-out test set reproduces
the same appreciable-band direction accuracy $99.8\%$ at detection-gate ROC-AUC $0.99$ and
calibration error $0.03$---so the accuracy is a property of the generative process, not of a
particular replicate draw.
It is also not limited by training-set size at this scale: doubling the training set to $6{,}400$
replicates, scored on the same held-out batch, moves rate-averaged accuracy by less than a tenth of a
point. The low- and mid-rate deficit is therefore information-limited rather than data-limited---the
empirical counterpart of the Fisher threshold (Corollary~\ref{cor:appreciable-threshold}), which caps
what any training-set size can recover at fixed genome size.
The discriminative content therefore lives in the rarefaction feature set itself, as the theory
predicts: the statistics are linear functionals of the joint site-frequency spectrum
(Theorem~\ref{thm:linear-sfs}), direction is identifiable from their first-order migration response
(Corollary~\ref{cor:identifiability}), and the across-locus variance supplies orientation even where
the mean cannot (Proposition~\ref{prop:variance-orients}). A linear (logistic) model on the same
features matches the network---on this benchmark it is marginally ahead, $78.0\%$ versus $75.6\%$
overall and essentially perfect in the appreciable band---exactly as
Proposition~\ref{prop:linear-optimal} predicts, since the direction geometry is first-order linear.
Both decisively exceed a precomputed Patterson's $D$, which is direction-blind by construction:
reversing the flow $j\!\to\!k$ leaves the symmetric allele-sharing pattern unchanged, so $D$ cannot
orient gene flow (Theorem~\ref{thm:symmetric}), and empirically it labels direction below the
majority-class rate. In positioning terms DNNaic is competitive among lightweight, outgroup-free,
summary-statistic direction methods---the lineage of ArchIE \citep{durvasula2019archie}---while
addressing a different task from genotype-matrix CNNs \citep{gower2021genomatnn,ray2024introunet}.
Migration-rate \emph{magnitude}, by contrast, remains information-limited at small rates, the
Fisher-information limit made precise in Theorem~\ref{thm:fisher}; the supported deliverable within
the stated simulation families is the direction of appreciable gene flow.

\paragraph{Which channels carry the orientation signal.} Restricting the classifier to feature
subsets under the same leakage-free protocol localizes the direction signal in the asymmetric,
population-labeled channels, exactly as Theorem~\ref{thm:symmetric} predicts
(Table~\ref{tab:channels}). The symmetric allelic-richness statistics alone orient poorly ($80.2\%$
in the appreciable band); the private and pairwise-private channels each do much better ($93.9\%$
and $98.9\%$) and together recover essentially the full-model accuracy ($99.6\%$ versus $99.8\%$).
The per-population labels that $D$-statistics discard are what make orientation possible.

\begin{table}[H]
\centering
\scriptsize
\begin{tabular}{lcc}
\toprule
feature channels & appreciable (\%) & overall (\%) \\
\midrule
allelic richness $\alpha$ only (symmetric) & 80.2 & 59.8 \\
private $\pi$ only & 93.9 & 67.9 \\
pairwise-private $\widehat\pi$ only & 98.9 & 73.4 \\
private $+$ pairwise-private (asymmetric) & 99.6 & 74.0 \\
all nine statistics (full) & \textbf{99.8} & \textbf{74.7} \\
\bottomrule
\end{tabular}
\caption{\textbf{Feature-channel ablation for direction} (HistGradientBoosting, leakage-free
$\mathrm{GroupKFold}(5)$, replicate-aggregated; appreciable $=$ rate $\ge 2.5\times10^{-4}$, chance
$33\%$). The asymmetric private and pairwise-private channels carry the orientation signal; the
symmetric richness statistics alone are far weaker, as Theorem~\ref{thm:symmetric} predicts.}
\label{tab:channels}
\end{table}

\paragraph{The across-locus variance orients beyond the mean.}
Proposition~\ref{prop:variance-orients} predicts that the across-locus \emph{variance} of the
rarefaction statistics carries orientation information the first-order \emph{mean} cannot, because
the rarefaction kernel weights a private allele linearly in the mean but quadratically in the
variance. A moment ablation under the identical leakage-free protocol confirms this directly
(Table~\ref{tab:moments}). Restricting the linear head to the across-locus \emph{means} of the nine
statistics gives $73.6\%$ rate-averaged direction accuracy; adding the \emph{variances} raises it to
$77.6\%$; and, most tellingly, the variances \emph{alone} ($75.2\%$) already exceed the means alone.
The second-moment coordinates are thus not redundant with the first---they supply independent
orientation signal, the mechanism Proposition~\ref{prop:variance-orients} isolates. (In the
appreciable band all four feature sets are near-saturated, so the discriminating axis is the
rate-averaged accuracy, which the low- and mid-rate replicates dominate.)

\begin{table}[H]
\centering
\scriptsize
\begin{tabular}{lcc}
\toprule
moments retained (linear head) & overall (\%) & appreciable (\%) \\
\midrule
across-locus mean only            & 73.6 & 98.9 \\
across-locus variance only        & 75.2 & 99.6 \\
mean $+$ variance                 & 77.6 & 100.0 \\
all three moments (mean, var, SE) & \textbf{78.1} & 100.0 \\
\bottomrule
\end{tabular}
\caption{\textbf{Moment ablation for direction} (logistic head, leakage-free
$\mathrm{GroupKFold}(5)$ on the $3{,}200$ replicates, replicate-aggregated, three-seed mean; chance
$33\%$). The across-locus variance alone out-orients the mean alone, and adding it to the mean
recovers essentially the full-feature accuracy---the empirical counterpart of
Proposition~\ref{prop:variance-orients}.}
\label{tab:moments}
\end{table}

\paragraph{Comparison with classical statistics.} The advantage over classical summaries is not
subtle. On a matched leakage-controlled benchmark, the \emph{same} learner given the rarefaction
features far outperforms it given the classical Patterson's $D$ and $f$-statistics
(Table~\ref{tab:sota}): a linear model reaches $91.5\%$ appreciable-regime three-way direction
accuracy on the rarefaction features versus $48.7\%$ on the $f$-statistics, and a precomputed
Patterson's-$D$ decision rule is degenerate, predicting the control for almost every replicate
($2.6\%$). The rarefaction features, not the model class, carry the orientation signal, and a linear
model on them is within a point of the network---exactly what Proposition~\ref{prop:linear-optimal}
leads one to expect.

\begin{table}[H]
\centering
\scriptsize
\begin{tabular}{lccc}
\toprule
method & overall (\%) & appreciable direction (\%) & macro-F1 \\
\midrule
Patterson's $D$ decision rule & 8.5 & 2.6 & 0.06 \\
$f$-statistics $+$ logistic & 39.2 & 48.7 & 0.31 \\
$f$-statistics $+$ gradient boosting & 36.5 & 39.3 & 0.30 \\
rarefaction $+$ logistic & 53.5 & \textbf{91.5} & 0.42 \\
rarefaction $+$ gradient boosting & 49.6 & 72.7 & 0.39 \\
rarefaction $+$ DNNaic (MLP) & 52.3 & 76.9 & \textbf{0.44} \\
\bottomrule
\end{tabular}
\caption{\textbf{Rarefaction features versus classical statistics}, all under the identical
leakage-controlled protocol on an independent $1{,}300$-replicate regeneration ($260$ held-out test
replicates; majority vote over rarefaction depths). Given the same learner, the ADZE rarefaction
features roughly double the appreciable-regime direction accuracy of the classical $f$-statistics,
and Patterson's $D$---direction-blind by construction (Theorem~\ref{thm:symmetric})---is degenerate.}
\label{tab:sota}
\end{table}

\paragraph{Cross-demography robustness.} To test whether the orientation signal is specific to the training
history, we froze the appreciable-direction model---fit on all $900$ in-distribution appreciable
replicates---and applied it, without any refitting, to fresh replicates simulated under alternative
demographies (a scoped probe of $45$ replicates per demography at rate $2.5\times10^{-4}$, on the
leakage-free rarefaction feature space that reproduces the in-distribution accuracy of $99.8\%$). Transfer is
essentially perfect to a fresh draw of the training demography and to a deeper, larger-$N_e$ history (both
$100\%$; $N_e=2\times10^{4}$, splits $1000/2000/4000$), and degrades gracefully to $73.3\%$ (bootstrap $95\%$
CI $[60,84]\%$, against a $33\%$ chance rate) under a recent, small-$N_e$ history ($N_e=5\times10^{3}$, splits
$250/500/1000$). The direction signal thus generalizes across demographies rather than being memorized from
the training history, with the expected attenuation when coalescence is fast and divergence shallow.

\subsection{A learned representation and a kernel correction do not beat the linear head}
\label{subsec:deep-kernel}

\begin{table}[b]
\centering
\begin{tabular}{llccc}
\toprule
Task & Metric & Linear & Network & Boosted trees \\
\midrule
Direction (classification) & accuracy (\%) & \textbf{78.9} & 76.0 & --- \\
Magnitude (regression) & MRE (\%), appreciable & 96.0 & 38.7 & \textbf{29.4} \\
Magnitude (regression) & Spearman, appreciable & 0.68 & 0.70 & \textbf{0.75} \\
\bottomrule
\end{tabular}
\caption{Linear model versus a network on the two halves of the problem, from one reproducible experiment
(\texttt{scripts/linear\_vs\_network.py}): one $54$-dimensional vector per replicate (mean and standard
deviation of the rarefaction coordinates across depths), leakage-free $5$-fold by replicate, fold-local
scaling. On direction the linear head is the most accurate---the first-order-linear classification, where
our theory predicts no classifier beats linear, and even a linear-skip or linear-initialised network
underperforms it. On magnitude the nonlinear models beat the linear ridge substantially (lower MRE is
better); gradient-boosted trees do as well as the network, so it is nonlinearity, not the network
specifically, that helps there.}
\label{tab:lin-vs-net}
\end{table}

Proposition~\ref{prop:linear-optimal} predicts that no classifier improves on the linear head to
first order in the migration magnitude, and the logistic-versus-network comparison of
Section~\ref{subsec:true-replicate} bears this out. That comparison leaves open whether a more
expressive model would do better, so we tested the strongest current recipe for extracting nonlinear
structure from a tabular problem: a neural network that learns a representation, followed by a
kernel-ridge correction fitted to the network's residuals on its own penultimate features. This is
the neural-mean-plus-kernel-correction construction used to build accurate physical-model emulators,
where the learned representation, not the kernel, is the dominant lever. Each replicate is summarized
by its rarefaction curve---the nine statistics at eight standardized depths together with their
across-depth mean; a four-layer residual network is trained on the training fold, and a
Mat\'ern-$5/2$ kernel ridge is fitted to its out-of-fold residuals, all under the same
$\mathrm{GroupKFold}(5)$ replicate split, so no replicate informs its own prediction at either stage.

The added machinery does not help. On three-way direction the network reaches $73.8\%$ overall
against the linear head's $79.0\%$, and the kernel correction leaves it essentially unchanged
($74.0\%$, fewer than one percent of the $2{,}700$ test labels changed); correcting the linear head
itself with the kernel holds it at $79.0\%$. The standard remedies for an underperforming network do not recover the linear head on direction either: in a companion experiment (\texttt{scripts/linear\_vs\_network.py}, Table~\ref{tab:lin-vs-net}), a residual network with an explicit linear skip warm-started from the logistic solution, a linear initialization, and heavy regularization all sit at $75$--$76\%$, short of the linear head. A linear skip guarantees the linear solution only if the nonlinear path is held at zero, but at this information limit that path has only noise to fit, and training spends its capacity on that noise at a cost to test accuracy. The robust choice for direction is therefore the linear head itself. On magnitude the network and gradient-boosted trees both beat a linear model but do no better than one another (Table~\ref{tab:lin-vs-net}), and the kernel correction does not improve the network. A positive control rules out an inert
pipeline: on synthetic data whose label is a nonlinear radius-band function of a rotated
two-dimensional subspace, the same network lifts accuracy from the linear head's $35\%$---the
three-way chance level---to $98\%$, and the kernel correction recovers a deliberately underfit
network's residual, from $81\%$ to $90\%$. The construction captures nonlinear structure and
structured residuals when they are present; on the rarefaction features it finds neither, because the
moment statistics are close to sufficient and the task sits at the Fisher-information limit of
Theorem~\ref{thm:fisher}. With the strongest available model this closes the possibility that the
linear head leaves accessible accuracy unused, and locates the leverage in the population-genetic
feature design rather than in classifier capacity. The comparison is reproduced by
\texttt{scripts/deep\_kernel\_comparison.py}.

\subsection{Scaling from three scenarios to all 12 ordered edges}
\label{subsec:twelve-directions}

The three-scenario benchmark asks whether a method can distinguish a small set of biologically
chosen alternatives. To test the harder class-count question directly, we constructed a separate
four-population benchmark containing every ordered donor--recipient edge
$\mathcal B_4=\{P_i\!\to P_j:i\ne j\}$, hence $K=12$. The baseline demography is deliberately
exchangeable: four equal-size populations ($N_e=10{,}000$) split simultaneously $500$
generations ago. One forward-time edge is active from the present to generation $250$ at
$m=10^{-3}$, giving integrated single-lineage exposure $mT=0.25$ and event probability
$1-e^{-0.25}=0.221$. Each of the $240$ independent replicates (20 per edge) contains a
$250$-kb sequence, 48 haploid gene copies per population, mutation rate $2\times10^{-8}$, and
recombination rate $1.78\times10^{-8}$. We audited the realized \texttt{msprime} epoch matrices
for all 12 edges: forward $P_i\!\to P_j$ is implemented as backward $P_j\!\to P_i$, and every
rate is exactly zero after generation $250$.

The four-population PADZE map contains four labelled richness coordinates $\alpha_i$, four
private-richness coordinates $\pi_i$, and all six pair-private coordinates
$\widehat\pi_{ij}$. For each channel we trace the across-locus mean, variance, and standard error
over rarefaction depths $g=2,\ldots,40$, then retain each curve's depth mean and standard
deviation, producing one 84-dimensional vector per independent simulation replicate. Evaluation uses
four-fold stratified cross-validation repeated at seeds 0, 1, and 2; no replicate crosses a fold.
For each replicate we majority-vote its three out-of-fold calls and report Wilson intervals over
the 240 independent simulation replicates. Exact-edge, unordered-pair, conditional orientation, donor, and
recipient chance rates are respectively $1/12$, $1/6$, $1/2$, $1/4$, and $1/4$.

The main intervention is dictated by Theorem~\ref{thm:many-directions}. In each training fold
only, we relabel every observation by all $|S_4|=24$ coordinate permutations and apply the same
permutation to its edge label; test observations are never augmented. This does not create new
independent data. It enforces the exact symmetry of the exchangeable baseline and reduces the
number of unrelated edge-specific rules that the classifier must estimate. The result is a
$12.9$-point gain in exact accuracy, from $120/240=50.0\%$ for ordinary logistic regression to
$151/240=62.9\%$ for its equivariant version (Table~\ref{tab:twelve-direction-models}); the three
single-repeat accuracies are $62.5\%$, $65.4\%$, and $65.4\%$. The equivariant rule recovers the
unordered pair at $76.3\%$ and, conditional on that pair being right, orients it at $82.5\%$.
The gain is therefore primarily in finding the interacting pair: ordinary logistic regression
already orients $81.1\%$ of the pairs it identifies. In a separate binary audit where the true
unordered pair is supplied, six independently fitted reversal classifiers orient $189/240=78.8\%$
(Wilson 95\% interval $[73.1,83.5]\%$); all six pairwise point estimates exceed chance, from
$70.0\%$ to $92.5\%$.

\begin{table}[H]
\centering
\scriptsize
\begin{tabular}{lccc}
\toprule
model & exact edge (\%, 95\% Wilson) & unordered pair (\%) & orientation $\mid$ pair (\%) \\
\midrule
logistic & 50.0 $[43.7,56.3]$ & 61.7 & 81.1 \\
\textbf{$S_4$-equivariant logistic} & \textbf{62.9 $[56.7,68.8]$} & \textbf{76.3} & \textbf{82.5} \\
Marchenko--Pastur bulk LDA & 48.8 $[42.5,55.0]$ & 64.6 & 75.5 \\
raw-feature RBF SVM & 49.2 $[42.9,55.5]$ & 62.5 & 78.7 \\
MLP & 45.4 $[39.2,51.7]$ & 57.1 & 79.6 \\
MLP features $+$ RBF SVM & 44.2 $[38.0,50.5]$ & 55.8 & 79.1 \\
MLP features $+$ MP whitening $+$ RBF SVM & 45.0 $[38.8,51.3]$ & 56.7 & 79.4 \\
population labels removed (control) & 7.1 $[4.5,11.1]$ & 14.2 & 50.0 \\
\bottomrule
\end{tabular}
\caption{\textbf{All 12 ordered directions on the exchangeable four-population baseline}
($n=240$ independent replicates; chance exact accuracy $8.3\%$). Entries are majority votes over
three repeated four-fold out-of-fold predictions. Permutation augmentation is confined to each
training fold. ``Orientation $\mid$ pair'' is evaluated only where the same multiclass call chose
the correct unordered pair; the oracle-pair reversal audit is reported in the text.}
\label{tab:twelve-direction-models}
\end{table}

Two controls locate the information. Sorting every exchangeable coordinate orbit removes the
population identities and returns $17/240=7.1\%$ exact accuracy; independently permuting replicate
labels before cross-validation gives $18/240=7.5\%$. Both agree with $1/12$ chance. Conversely,
added capacity supplies no hidden gain. A raw RBF kernel, an MLP, and an RBF kernel on learned MLP
features all fall below the ordinary linear head. We also tested a random-matrix correction
\citep{marchenko1967distribution}: within each training fold, pooled within-class covariance
eigenvalues inside the estimated Marchenko--Pastur bulk are collapsed to the noise level while
outlying spikes are retained. This MP-bulk LDA reaches $48.8\%$; applying the corresponding
whitener to the MLP representation before an RBF kernel reaches $45.0\%$. The neural/RMT hybrid is
therefore a negative result, not a post hoc win: at this sample size the useful inductive bias is
known population symmetry, not additional nonlinearity.

We finally froze every baseline model and changed one factor at a time, with 96 fresh replicates
per family (eight per edge). Table~\ref{tab:twelve-direction-transfer} separates frozen transfer
from retained information by also refitting the equivariant logistic rule within each shifted
family under the same leakage-free folds. Fivefold weaker exposure approaches chance even after
refitting, so the limiting factor is signal. A fixed $N_e$ imbalance and a labelled balanced
topology reduce frozen accuracy; their partial or full recovery after refitting diagnoses
distribution shift. Halving sequence length lowers both frozen and refitted accuracy to about
$50\%$. These small shifted panels are stress tests, not precision benchmarks, but they reject a
blanket transfer claim. In particular, the risk guarantee in Theorem~\ref{thm:many-directions}
requires an $S_s$-invariant base law; the fixed labelled balanced topology intentionally violates
that premise.

\begin{table}[H]
\centering
\scriptsize
\begin{tabular}{lcc}
\toprule
test family & frozen equivariant (\%, 95\% Wilson) & refit within family (\%, 95\% Wilson) \\
\midrule
weak signal, $mT=0.05$ & 16.7 $[10.5,25.4]$ & 10.4 $[5.8,18.1]$ \\
unequal $N_e=(5,10,20,8)\times10^3$ & 45.8 $[36.2,55.8]$ & 56.3 $[46.3,65.7]$ \\
balanced $((P_1,P_2),(P_3,P_4))$ tree & 35.4 $[26.6,45.4]$ & 64.6 $[54.6,73.4]$ \\
half sequence, 125 kb & 49.0 $[39.2,58.8]$ & 51.0 $[41.2,60.8]$ \\
\bottomrule
\end{tabular}
\caption{\textbf{Frozen nuisance transfer for the 12-edge classifier} ($n=96$ independent
replicates per row; exact-edge chance $8.3\%$). The frozen column applies the baseline model once,
without refitting. The last column is a fresh repeated four-fold cross-validation within the
shifted family and is included only to distinguish loss of transfer from loss of information.}
\label{tab:twelve-direction-transfer}
\end{table}

The complete seed ledger, realized migration epochs, fold-overlap audit, classwise confusion
matrices, covariance-spectrum diagnostics, and exact artifacts are emitted by the CPU-only,
checkpointed script \texttt{twelve\_direction\_extension.py}.

\subsection{Real-data application: 1000 Genomes}
\label{subsec:real-data}

To exercise the method end-to-end on empirical data, we ran the released \textsc{PADZE} directly on
1000 Genomes Project VCFs \citep{tgp2015global}, fetching the same five genome-spanning windows used
to validate the tool and computing the $28$-dimensional feature vector for several three-population
trios. The extraction reproduces the benchmark exactly---$9{,}756$ polymorphic biallelic SNPs on the
YRI/CEU/CHB trio---so \textsc{PADZE} runs on real project data with no bespoke tooling.

The feature computation is faithful, and the real vectors land close to the simulation prior.
Allelic-richness moments recover the expected diversity ordering and topology: on YRI/CEU/CHB the
out-of-Africa gradient appears (YRI $1.669 >$ CEU $1.423 \approx$ CHB $1.421$ at depth $g=100$), the
two sister populations show near-equal richness, and the pairwise-private structure gives the correct
tree. The resulting feature vectors sit largely \emph{within} the simulated range (median per-feature
deviation about $0.7$ standard deviations, with a few coordinates in the tails), so the sim-trained
model applied without refitting is interpolating on most features rather than extrapolating wildly.

\paragraph{Validated recovery on real genotype backgrounds.} No natural human trio carries a
labelled introgression event, so we validated direction recovery directly on real genetic data by
\emph{injecting} a known signal. Starting from the CEU/CHB/YRI trio (CEU and CHB the sister pair,
YRI the outgroup), whose gate abstains in its natural state, we turned a fraction $f$ of a chosen
recipient population's individuals into migrants carrying a donor individual's genotypes---a
controlled analogue of directional gene flow imposed on real allele frequencies and linkage---and
ran the released \textsc{PADZE} and the simulation-trained model on the result (four replicates per
condition). On these real backgrounds the method behaves as the theory predicts. It abstains on the
un-injected control (mean gate probability $0.01$), and its detection score rises with the injected
fraction ($0.34$, $0.52$, $0.70$ at $f=0.05,0.1,0.2$). Once the injected signal is appreciable
($f\ge0.1$), three-way direction accuracy is $100\%$ ($16/16$ over the two rates and four
replicates)---including the difficult reversed pair $B$ (CHB$\to$YRI) versus $C$ (YRI$\to$CHB)---and
at a weak injection ($f=0.05$) it falls to $62\%$, the same appreciable-rate dependence seen in the
simulations (Section~\ref{subsec:true-replicate}). The direction signal the method reads is thus
recoverable on real human variation, not an artifact of the coalescent simulator.

On the natural trios we make no accuracy claim---they carry no labelled within-population introgression,
and rarefaction statistics are demography-dependent (Theorem~\ref{thm:expansion})---so every
real-data call is a behavioural demonstration, not a measurement. What the demonstrations show is
that the detect-then-orient gate tracks the relatedness of the sampled populations in the expected
direction. Across three trios (all with in-distribution features, richness $1.42$--$1.71$), it
abstains on the deeply differentiated YRI/CEU/CHB trio (appreciable-migration probability $0.01$),
where the three populations are cleanly separated and any gene flow is ancient, and it fires
($0.90$--$0.93$) on the two trios built around a recently diverged sister pair---two Europeans
(GBR/TSI) or two Han Chinese (CHB/CHS), each with an African outgroup---where the sisters share
extensive variation. This is the logic the method encodes, high inter-population sharing reading as
gene flow; but on human trios recent shared ancestry and ongoing gene flow are confounded, so we
cannot call these detections right or wrong. A useful natural-data test must therefore combine
independent biological context with a negative control under one fixed sampling and feature
protocol. The next analysis shows why the control is load-bearing.

\subsection{Real-data stress test: signal and specificity in \textit{Heliconius}}
\label{subsec:heliconius}

The \textit{Heliconius melpomene}--\textit{timareta}--\textit{cydno} complex is a useful stress
test because adaptive transfer of red wing-pattern alleles from \textit{H.~melpomene} into
\textit{H.~timareta} is documented at particular loci
\citep{pardodiaz2012adaptive,wallbank2016evolutionary,martin2013speciation}. On the initial
genome-wide trio, Patterson's $D$ detects excess sharing,
\[
D(\textit{cydno},\textit{timareta},\textit{melpomene};\textit{numata})=0.51
\qquad (Z=33),
\]
and the raw pair-private ratio is
\[
\frac{\widehat\pi_{\textit{timareta},\textit{melpomene}}}
     {\widehat\pi_{\textit{cydno},\textit{melpomene}}}\simeq 2.6.
\]
The frozen simulation-trained
direction head also calls class $C$ ($\textit{melpomene}\!\to\!\textit{timareta}$). Those three
agreements motivated the original application, but one positive panel cannot establish natural-data
specificity.

We therefore reran the analysis as a fixed-protocol control audit. Each panel uses exactly ten
diploid individuals per focal population, requires at least 16 called gene copies in every
population, and retains a seeded reservoir of $15{,}000$ SNPs over all 21 chromosomes. Four
geographically motivated panels have raw pair-private ratios $1.882$--$2.563$, and the primary
$54$-dimensional linear head calls $C$ on each with uncalibrated softmax scores
$0.943$ to above $0.999999$. Crucially, the intended allopatric \textit{mel\_ros} control behaves
differently in the data but not in the head: its earlier $D$ estimate is approximately zero
($D=0.01$), its raw pair-private ratio reverses to $0.626$, yet the head still calls $C$ with score
$0.9918$. Thus Patterson's $D$ and the model-free sharing contrast distinguish the selected
geographic signal from this control, whereas the frozen orientation head has no specificity on
these natural panels.

The same distinction changes how stability must be read. All 21 leave-one-chromosome-out versions
of the first panel call $C$ at scores $0.999999982$--$0.999999991$, but a stable extrapolation that
also assigns $C$ to the control is evidence of systematic simulation-to-natural miscalibration, not
external validation. The panels reuse populations, individuals, and linked loci; their exact
genome-wide donor--recipient labels are not independently known; and adaptive-locus direction need
not label a whole-genome mixture. We therefore retain the biological findings---$D$ detects excess
sharing and population-labelled rarefaction exposes an interpretable asymmetry---but treat the
trained natural-panel direction scores as an exploratory failure mode. Target-demography training
and genuinely independent labelled controls are required before a natural direction call is
calibrated. The full sample, filter, chromosome, source-hash, score, and sensitivity ledger is in
\texttt{results/natural\_diagnostics\_2026\_07\_09/heliconius.json}.

\subsection{Real-data application: Neanderthal introgression into non-Africans}
\label{subsec:neanderthal}

The most thoroughly documented introgression in humans is the Neanderthal ancestry carried by
non-Africans---on the order of two percent, and essentially absent from sub-Saharan Africans---the
result that first established the ABBA--BABA test \citep{green2010neandertal,prufer2014complete}. We
ran the released \textsc{PADZE} directly on real genotypes for this system, slicing five
genome-spanning windows from the 1000 Genomes phase-3 VCFs \citep{tgp2015global} and from the two
high-coverage Neanderthal genomes---Altai and Vindija~33.19
\citep{prufer2014complete,prufer2017vindija}---on the same GRCh37 coordinates, and intersecting to
$21{,}702$ biallelic SNPs. The trio maps to the caterpillar tree as $P_1=$~African (YRI),
$P_2=$~Eurasian, $P_3=$~Neanderthal, so Neanderthal$\to$Eurasian is the class-$C$ orientation.

The detection reproduces cleanly, with the population contrast it should have. Patterson's statistic is
\[
D(\text{YRI},\text{CEU},\text{Neanderthal};\text{ancestral})=+0.27
\qquad (Z=3.4;\ \mathrm{ABBA}=76,\ \mathrm{BABA}=44).
\]
Thus the European shares an excess of derived alleles with
the Neanderthal that the African does not, the signature of Neanderthal$\to$non-African gene flow,
and the East-Asian trio is likewise positive ($D=+0.20$). The excess collapses in both controls,
exactly as it must:
\[
D(\text{YRI},\text{LWK},\text{Neanderthal})=+0.03,
\qquad
D(\text{CEU},\text{CHB},\text{Neanderthal})=-0.06.
\]
The first has ABBA $\approx$ BABA; the second is the near-zero and faintly
East-Asian-leaning value expected when both carry comparable Neanderthal ancestry. The labelled
ABBA/BABA asymmetry localizes the excess sharing to the Eurasian rather than the African---the
per-population information a symmetric $D$ discards---while $D$ itself, being symmetric, still cannot
name the donor.

This system also marks the operating boundary of the rarefaction route to \emph{orientation}, and it
is a boundary the method respects on its own. The learned gate abstains on all four Neanderthal trios
(mean appreciable probability $0.00$) rather than asserting a direction, and two compounding causes
explain why. First, only a handful of high-coverage archaic genomes exist, so the Neanderthal
population supplies four gene copies and the common rarefaction depth is $g\le 4$; on the simulations,
restricting every population to $g\le 4$ lowers appreciable-band direction accuracy from $99.9\%$ to
$84.1\%$ and the detection gate from ROC-AUC $0.99$ to $0.86$ (Fisher linear discriminant,
leakage-free)---degraded but not destroyed, so shallow depth alone does not force the abstention.
Second, and more decisively, the Neanderthal lineage is far more diverged than any split in the
training demography, so its rarefaction features fall well outside the simulated range (median
per-feature deviation about three standard deviations); the calibrated gate responds to
out-of-distribution input by abstaining rather than extrapolating---the safe behaviour the two-stage
design intends. The site-based Patterson's $D$, which does not rarefy and assumes no training
distribution, is unaffected and still detects. Sample depth is therefore necessary but not sufficient
for transfer. The Neanderthal gate abstains on shallow, deeply divergent input, whereas the
\textit{Heliconius} head extrapolates confidently even on an allopatric control
(Section~\ref{subsec:heliconius}); the learned system is not a universal out-of-distribution guard.
Population-scale, target-matched simulations and explicit controls are needed before orientation,
while few-genome archaic introgression remains detectable by $D$ but not orientable by these frozen
features.

\section{Theoretical results connecting rarefaction statistics and evaluation}
\label{sec:theoretical-results}

This section collects the mathematical results that make the methodology more than an
empirical neural-network fit; each statement is proved inline.

\begin{thm}[Rarefaction statistics are linear SFS functionals]\label{thm:linear-sfs}
Fix sample sizes and rarefaction depth \(g\). Let \(\Xi_{a,b,c}\) denote the joint sample
frequency spectrum count for alleles observed \(a,b,c\) times in the three sampled
populations. For every rarefaction statistic
\[
T_g\in\{\alpha_1,\alpha_2,\alpha_3,\pi_1,\pi_2,\pi_3,
\widehat{\pi}_{12},\widehat{\pi}_{13},\widehat{\pi}_{23}\},
\]
there is a deterministic combinatorial weight vector \(w_{T,g}\) such that
\[
\mathbb{E}[T_g]=\langle w_{T,g},\mathbb{E}[\Xi]\rangle .
\]
\end{thm}

\begin{proof}
Condition on a site with derived-allele counts \((a,b,c)\). After rarefying population
\(j\) to depth \(g\), the probability that the allele is present is a hypergeometric
probability depending only on the original count and sample size. Private and
pair-private events are finite Boolean combinations of these rarefied-presence events. For
example, the event counted by \(\pi_1\) is ``present after rarefaction in population 1 and
absent after rarefaction in populations 2 and 3.'' Therefore the conditional contribution
of a site with counts \((a,b,c)\) to \(T_g\) is a deterministic weight
\(w_{T,g}(a,b,c)\). Summing these weights against the joint SFS gives \(T_g\), and taking
expectations gives the displayed linear functional.
\end{proof}

\begin{thm}[Small-migration expansion]\label{thm:expansion}
Consider a structured coalescent model whose generator can be written
\[
Q_\beta(m)=Q_0+mQ_\beta
\]
for direction \(\beta\) and migration magnitude \(m\). For every bounded rarefaction
statistic \(T_g\),
\[
\mathbb{E}_{\beta,m}[T_g]=\mathbb{E}_{0}[T_g]+m\Delta_{\beta,g,T}+O(m^2)
\]
as \(m\to 0\). The first-order coefficient \(\Delta_{\beta,g,T}\) is obtained by inserting
one \(Q_\beta\) migration transition into the no-migration ancestral semigroup.
\end{thm}

\begin{proof}
For fixed sample size and a finite state truncation of the ancestral process, expectations
of bounded genealogical observables are finite sums of transition probabilities from the
semigroup \(\exp(t(Q_0+mQ_\beta))\). Differentiating the semigroup at \(m=0\) gives
\[
\frac{d}{dm}\exp(t(Q_0+mQ_\beta))\bigg|_{m=0}
=\int_0^t \exp(sQ_0)Q_\beta\exp((t-s)Q_0)\,ds .
\]
This derivative corresponds to histories with exactly one migration event of type
\(\beta\). Mutation and the map from genealogies to the expected joint SFS are linear, and
Theorem 1 shows that rarefaction statistics are linear SFS functionals. Composing these
linear maps with the differentiated semigroup gives the coefficient
\(\Delta_{\beta,g,T}\), and analyticity of finite-dimensional matrix exponentials gives
the \(O(m^2)\) remainder.
\end{proof}

\begin{cor}[Generic local identifiability of direction]\label{cor:identifiability}
Let
\[
\Delta_\beta=(\Delta_{\beta,g,T})_{g,T}
\]
be the vector of first-order perturbations over the rarefaction depths and statistics
used by the model. If
\[
\min_{\beta\ne\beta'}\|\Delta_\beta-\Delta_{\beta'}\| \ge \gamma>0,
\]
then there exists \(m_0>0\) such that, for all \(0<m<m_0\),
\[
\|\mathbb{E}_{\beta,m}[\Phi]-\mathbb{E}_{\beta',m}[\Phi]\|
\ge \frac{\gamma}{2}m .
\]
Thus introgression direction is locally identifiable from the expected rarefaction feature
vector \(\Phi\).
\end{cor}

\begin{proof}
Subtract the expansions in Theorem 2 for two directions:
\[
\mathbb{E}_{\beta,m}[\Phi]-\mathbb{E}_{\beta',m}[\Phi]
=m(\Delta_\beta-\Delta_{\beta'})+O(m^2).
\]
The assumed separation gives a first-order norm at least \(\gamma m\). For sufficiently
small \(m\), the quadratic remainder has norm at most \(\gamma m/2\), proving the stated
lower bound.
\end{proof}

The next theorem is the complementary impossibility result: it explains \emph{why} the
asymmetric private and pair-private rarefaction components are necessary, by showing that the
symmetric, \emph{population-unlabeled} summaries underlying classical pairwise tests cannot
orient gene flow at all. This is the theoretical counterpart of two empirical observations: in
PC space the recently diverged pair $B$ (P2$\to$P3) and $C$ (P3$\to$P2) overlaps most
(Section~\ref{sec:eda}, Figure~\ref{fig:pca-comparison}), and under replicate-aware evaluation this donor--recipient reversed pair shows elevated mutual confusion. The escape route the rarefaction features take is to retain
the per-population labeling that $D$-statistics discard.

\begin{thm}[Symmetric outgroup-free summaries cannot orient gene flow]\label{thm:symmetric}
Suppose populations $j$ and $k$ are exchangeable under the no-migration demography (equal
effective sizes and a divergence history symmetric in $j\leftrightarrow k$), and let $S$ be
any statistic that depends on the genealogy only through symmetric pairwise allele-sharing
summaries $\{s_{jk}\}$---without an outgroup, without the temporal order of sampling, and
without asymmetric private or shared-private components. Let history $H$ carry directional
gene flow $j\to k$, and let $H'=\tau\!\cdot\! H$ be its image under the label transposition
$\tau=(j\,k)$, so that $H'$ carries gene flow $k\to j$. Then
\[
\mathrm{Law}(S\mid H)=\mathrm{Law}(S\mid H'),
\]
and the direction of gene flow is not identifiable from $S$. The private and pair-private
rarefaction statistics $\pi_j,\widehat\pi_{jk}$ used by DNNaic are not $\tau$-invariant and
therefore escape this obstruction.
\end{thm}

\begin{proof}
Because $j$ and $k$ are exchangeable under the no-migration demography, the transposition
$\tau=(j\,k)$ is a measure-preserving symmetry of the ancestral process \emph{except} for the
directional migration edge, whose orientation it reverses; hence $\tau\!\cdot\! H=H'$ as laws
on genealogies. Any statistic built from the symmetric summaries $\{s_{jk}\}$ satisfies
$S\circ\tau=S$, so
\[
\mathrm{Law}(S\mid H')=\mathrm{Law}(S\circ\tau\mid H)=\mathrm{Law}(S\mid H).
\]
The two histories thus share the law of $S$ but differ in direction, which is exactly
non-identifiability of direction from $S$. Finally, $\pi_j$ counts alleles private to
population $j$ and $\widehat\pi_{jk}$ counts alleles shared by $j,k$ but absent from the third
population; under $\tau$ these obey $\pi_j\leftrightarrow\pi_k$ and
$\widehat\pi_{jk}\leftrightarrow\widehat\pi_{jk}$, and once migration breaks exchangeability
$\pi_j\neq\pi_k$ generically, so the private/pair-private channel retains orientation
information that the symmetric summaries discard.
\end{proof}

\paragraph{Second moments orient gene flow where the first-order mean is blind.}
Corollary~\ref{cor:identifiability} recovers direction from the \emph{mean} rarefaction features
whenever their first-order migration response separates two directions. Because the rarefaction summaries also
report the across-locus \emph{variance} of every statistic, it is natural to ask whether that
second-moment coordinate adds orientation information beyond the mean. The next proposition---a
special case of a moment-hierarchy phenomenon---answers yes: it exhibits reversed directions
whose first-order \emph{mean} margin vanishes while the first-order \emph{variance} margin does
not. It thereby strengthens Corollary~\ref{cor:identifiability} by supplying an identifying
coordinate exactly in the degenerate regime where that corollary's mean-based margin $\gamma$ is
zero.

We work in the standard independent-loci model \citep{sawyer1992population}: unlinked biallelic
loci are i.i.d.\ draws from the genealogy law $F_{\beta,m}$, and the across-locus mean and
variance of a statistic are the sample mean and variance of its per-locus values, hence
consistent estimators of the per-locus mean and variance under $F_{\beta,m}$. For a derived
allele with sample configuration $(a,b,c)$---derived counts in the exchangeable pair $(j,k)$ and
the third population $\ell$---write the depth-$g$ rarefied presence/absence probabilities as
$p^{(N)}_a=1-Q(a,g;N)$ and $q^{(N)}_a=Q(a,g;N)$. Assuming, as in the private-allele regime, that
the ancestral allele is rarefied-present in all three populations (so only the derived allele
can be private), the per-locus private-$j$ richness at depth $g$ is the bounded configuration
weight
\[
\pi_j(g)=w_j(a,b,c;g),\qquad w_j(a,b,c;g)=p^{(N_j)}_a\,q^{(N_k)}_b\,q^{(N_\ell)}_c,
\]
of Theorem~\ref{thm:linear-sfs}. Writing $f_{a,b,c}$ for the per-locus configuration law and
$\langle f,h\rangle=\sum_{a,b,c}f_{a,b,c}\,h(a,b,c)$, the rarefaction mean and variance estimate
$\mu_j(g)=\langle f,w_j\rangle$ and $v_j(g)=\langle f,w_j^2\rangle-\langle f,w_j\rangle^2$. We
adopt the exchangeable setup of Theorem~\ref{thm:symmetric}: the no-migration demography is
symmetric under $\tau=(j\,k)$ (equal sample sizes $N_j=N_k=N$ and a
$j\!\leftrightarrow\!k$-symmetric divergence history), the reversed directions
$\beta=(j\!\to\!k)$ and $\beta'=\tau\!\cdot\!\beta=(k\!\to\!j)$ obey
$f^{\beta'}_{a,b,c}=f^{\beta}_{b,a,c}$, and by Theorem~\ref{thm:expansion} each law is analytic,
$f^{\beta}=f^{0}+m\,\delta^{\beta}+O(m^2)$ with $\langle\delta^\beta,\mathbf 1\rangle=0$ and
$\tau$-invariant base $f^0_{a,b,c}=f^0_{b,a,c}$.

\begin{prop}[Variance orients gene flow at first order where the mean does not]
\label{prop:variance-orients}
Define at depth $g$ the first-order mean and variance \emph{margins} between the reversed
directions,
\[
M_\mu(g)=\partial_m\!\left[\mu_j^{\beta}(g)-\mu_j^{\beta'}(g)\right]_{m=0},\qquad
M_v(g)=\partial_m\!\left[v_j^{\beta}(g)-v_j^{\beta'}(g)\right]_{m=0}.
\]
Then, with $u_g:=w_j-w_j\!\circ\!\tau$, $r_g:=w_j^2-w_j^2\!\circ\!\tau$ and
$\mu^0(g)=\langle f^0,w_j\rangle$,
\[
M_\mu(g)=\langle\delta^{\beta},u_g\rangle,\qquad
M_v(g)=\langle\delta^{\beta},r_g\rangle-2\mu^0(g)\,M_\mu(g),
\]
and, writing $q_c:=q^{(N_\ell)}_c$,
\[
u_g(a,b,c)=q_c\,(p^{(N)}_a-p^{(N)}_b),\qquad
r_g(a,b,c)=q_c^{\,2}\,(p^{(N)}_a-p^{(N)}_b)\big(p^{(N)}_a q^{(N)}_b+q^{(N)}_a p^{(N)}_b\big).
\]
The kernels $u_g$ and $r_g$ are linearly independent, so $M_\mu$ and $M_v$ are independent
linear functionals of the migration response $\delta^\beta$. Consequently there exist responses
that are mean-blind but variance-visible, $M_\mu(g)=0\neq M_v(g)$. One such response lives on a
singleton and a higher-count private class: with $c=0$, count $h$ ($1<h\le N$),
$p_h:=1-Q(h,g;N)$, and pair differentials
$\eta:=\delta^\beta(h,0,0)-\delta^\beta(0,h,0)$,
$\xi:=\delta^\beta(1,0,0)-\delta^\beta(0,1,0)$, imposing $M_\mu(g)=0$ fixes
$\xi=-\tfrac{Np_h}{g}\,\eta$ and yields
\[
M_v(g)=p_h\!\left(p_h-\tfrac{g}{N}\right)\eta,\qquad
\big|M_v(g)\big|=p_h\!\left(p_h-\tfrac{g}{N}\right)|\eta|>0\ \ (\eta\neq0),
\]
the factor $p_h-g/N>0$ because $Q(h,g;N)<Q(1,g;N)$ for $h>1$. The reversed directions are thus
first-order indistinguishable from the mean private-richness coordinate yet separated at first
order by its across-locus variance.
\end{prop}

\begin{proof}
As $\pi_j(g)=w_j$ is a bounded configuration function, the across-locus sample mean and variance
converge to $\mu_j=\langle f,w_j\rangle$ and $v_j=\langle f,w_j^2\rangle-\langle f,w_j\rangle^2$;
analyticity of $f^\beta$ in $m$ (Theorem~\ref{thm:expansion}) gives
$\partial_m\mu_j^\beta|_0=\langle\delta^\beta,w_j\rangle$ and
$\partial_m v_j^\beta|_0=\langle\delta^\beta,w_j^2\rangle-2\langle f^0,w_j\rangle\langle\delta^\beta,w_j\rangle$.
Let $\tau$ act on configurations by $(a,b,c)\mapsto(b,a,c)$; since $N_j=N_k=N$ we have
$w_j\!\circ\!\tau=w_k$, and $\delta^{\beta'}=\delta^{\beta}\!\circ\!\tau$ with the involution
identity $\langle g\!\circ\!\tau,h\rangle=\langle g,h\!\circ\!\tau\rangle$ gives
$\partial_m\mu_j^{\beta'}|_0=\langle\delta^\beta,w_k\rangle$. Hence
$M_\mu=\langle\delta^\beta,w_j-w_k\rangle=\langle\delta^\beta,u_g\rangle$. The same substitution
in the second moment, with $f^0\!\circ\!\tau=f^0$ so that
$\langle f^0,w_j\rangle=\langle f^0,w_k\rangle=\mu^0$, gives
$M_v=\langle\delta^\beta,w_j^2-w_k^2\rangle-2\mu^0\langle\delta^\beta,u_g\rangle
=\langle\delta^\beta,r_g\rangle-2\mu^0 M_\mu$.
Writing $p_a=p^{(N)}_a$, $q_a=1-p_a$ and $q_c=q^{(N_\ell)}_c$, we have $w_j=p_aq_bq_c$,
$w_k=q_ap_bq_c$, so $u_g=q_c(p_aq_b-q_ap_b)=q_c(p_a-p_b)$ and
$r_g=(w_j-w_k)(w_j+w_k)=q_c^2(p_a-p_b)(p_aq_b+q_ap_b)$; their ratio
$r_g/u_g=q_c(p_aq_b+q_ap_b)$ is non-constant, so $u_g,r_g$ are not proportional and
$\{\langle\delta,u_g\rangle=0\}\not\subseteq\{\langle\delta,r_g\rangle=0\}$. For the witness,
restrict $\delta^\beta$ to $(1,0,0),(h,0,0),(0,1,0),(0,h,0)$ with a mass-conserving remainder on
$\tau$-fixed configurations where $w_j=w_k=0$. There $q_c=1$ and, by the exact singleton
identity $p_1=1-Q(1,g;N)=g/N$,
\[
M_\mu=\tfrac{g}{N}\,\xi+p_h\,\eta,\qquad
\langle\delta^\beta,r_g\rangle=\big(\tfrac{g}{N}\big)^2\xi+p_h^2\,\eta .
\]
Setting $M_\mu=0$ gives $\xi=-\tfrac{Np_h}{g}\eta$, whence (using $M_\mu=0$ once more)
$M_v=\langle\delta^\beta,r_g\rangle=\big(\tfrac{g}{N}\big)^2\big(\!-\tfrac{Np_h}{g}\eta\big)+p_h^2\eta
=p_h\big(p_h-\tfrac{g}{N}\big)\eta$. Since $h>1$ implies $Q(h,g;N)<Q(1,g;N)$, i.e.\ $p_h>g/N$,
this is nonzero for $\eta\neq0$. Finally $f^0+m\,\delta^\beta$ is a probability law for small
$|m|$ (the four base masses are positive), and the reversed-direction law is its $\tau$-image, as
the model requires; the witness is therefore admissible.
\end{proof}

\begin{rem}
Within the independent-loci model the margin identities and the explicit witness are exact, so
Proposition~\ref{prop:variance-orients} is proved as an existence statement. It does \emph{not}
assert that a particular named demography realizes the exact cancellation $M_\mu(g)=0$: whether a
given structured-coalescent demography places $\delta^\beta$ on the mean-blind hyperplane is a
single (codimension-one) condition on the demographic parameters, checkable by finite differences
in \texttt{msprime} and left to the numerical program noted in Section~\ref{sec:limitations}. The
unconditional content is what matters for the method---first-order mean-blindness does not imply
first-order variance-blindness, because the rarefaction kernel weights a private allele
\emph{linearly} in the mean ($w_j$) but \emph{quadratically} in the variance ($w_j^2$), so
equalizing the expected private count across reversed directions leaves the across-locus
concentration unequal. This is the mechanism by which the second-moment coordinates of the rarefaction summaries
carry orientation beyond the first moment, matching the measured importance of the
variance-bearing private and pair-private channels. The moment ablation of
Table~\ref{tab:moments} confirms the prediction directly: under the leakage-free protocol the
across-locus variance alone out-orients the mean alone, and adding it to the mean recovers the
full-feature accuracy.
\end{rem}

The identifiability results above separate the directions through the first-order response
vectors $\Delta_\beta$, but they do not say what \emph{kind} of classifier is needed to exploit
that separation. The next proposition answers this, and in doing so explains why a plain linear
model on the rarefaction features performs on par with the deep network: in the small-migration
regime the discriminative geometry is linear, so a linear rule is already optimal.

\begin{prop}[A linear rule is first-order optimal for direction]\label{prop:linear-optimal}
Work in the local asymptotic normality regime of Theorem~\ref{thm:fisher}: the per-replicate
rarefaction feature vector $\Phi$ obeys $\sqrt{L}\,(\Phi-\mathbb{E}_{\beta,m}\Phi)\to
\mathcal{N}(0,\Sigma_0)$ as the number of approximately independent loci $L\to\infty$, with a
common leading covariance $\Sigma_0\succ0$ that does not depend on the direction $\beta$ because
the no-migration base demography is exchangeable in the experimental populations. By
Theorem~\ref{thm:expansion} the class-conditional mean is
$\mu_\beta(m)=\mu_0+m\,\Delta_\beta+O(m^2)$. Write $Q(z)=\Pr(\mathcal N(0,1)>z)$. Then, to first
order in $m$,
\begin{enumerate}
\item the Bayes-optimal rule separating any two directions $\beta,\beta'$ is the \emph{linear}
discriminant
\[
\phi\ \longmapsto\ \operatorname{sign}
\big\langle \Sigma_0^{-1}(\Delta_\beta-\Delta_{\beta'}),\ \phi-\tfrac12(\mu_\beta+\mu_{\beta'})\big\rangle;
\]
\item its error is the Gaussian overlap
$Q\!\big(\tfrac12\, m\sqrt{L}\,\lVert\Sigma_0^{-1/2}(\Delta_\beta-\Delta_{\beta'})\rVert\big)$,
which decreases monotonically in $m$ and tends to the two-class chance level $\tfrac12$ as
$m\to0$;
\item consequently no classifier---linear or nonlinear---attains a smaller first-order error than
a linear model on the rarefaction features.
\end{enumerate}
\end{prop}

\begin{proof}
Under the stated law the per-replicate feature is, to leading order, Gaussian with
direction-independent covariance $\Sigma_0/L$ and mean $\mu_\beta(m)$. For two Gaussians with a
common covariance the log-likelihood ratio is affine in $\phi$, so the Bayes rule is the
hyperplane in (i) (Fisher's linear discriminant), and the two-class Bayes error equals
$Q(\Delta_{\mathrm M}/2)$ with Mahalanobis separation
$\Delta_{\mathrm M}=\lVert(\Sigma_0/L)^{-1/2}(\mu_\beta-\mu_{\beta'})\rVert
=m\sqrt{L}\,\lVert\Sigma_0^{-1/2}(\Delta_\beta-\Delta_{\beta'})\rVert+O(m^2)$, which is (ii). Because
the optimal boundary is already affine, the linear hypothesis class contains the Bayes rule, so
enlarging the class cannot lower the error to this order, giving (iii).
\end{proof}

\begin{rem}
Proposition~\ref{prop:linear-optimal} accounts for two measured facts at once. First, a logistic
model on the rarefaction features matches the deep network to within about a point of accuracy
(Section~\ref{subsec:true-replicate}): the direction geometry is first-order linear, so the
network's nonlinearity is not what carries the signal---the population-labeled private and
pair-private channels are. Second, the error formula $Q(\tfrac12 m\sqrt L\lVert\cdots\rVert)$
reproduces the shape of the per-rate accuracy curve (Table~\ref{tab:true-replicate}): it collapses
to chance as the migration magnitude falls below a Fisher threshold of the same $L^{-1/2}$ order as
the magnitude threshold $m_\star$ of Corollary~\ref{cor:appreciable-threshold}. The two thresholds
are set by \emph{different} functionals of the same geometry, however---magnitude by the individual
shift $\lVert\Sigma_0^{-1/2}\Delta_\beta\rVert$ and orientation by the pairwise differences
$\lVert\Sigma_0^{-1/2}(\Delta_\beta-\Delta_{\beta'})\rVert$---as made precise in
Proposition~\ref{prop:two-thresholds}. The small residual advantage of the network at larger
rates is consistent with the neglected $O(m^2)$ and any direction-dependent curvature of $\Sigma$,
which a nonlinear model can use once the migration signal is strong.
\end{rem}

The final perturbative result quantifies the empirical difficulty of magnitude estimation: it
turns the weak small-rate relative error into a principled information limit rather than a
modeling failure.

\begin{thm}[Local rate estimability and Fisher information]\label{thm:fisher}
Fix a direction $\beta$ with nonzero first-order response $\Delta_\beta\neq 0$ from
Theorem~\ref{thm:expansion}. Assume the per-locus empirical rarefaction feature vector obeys a
local asymptotic normality law $\sqrt{L}\,(\widehat\Phi-\mathbb{E}_m\Phi)\to
\mathcal{N}(0,\Sigma_0)$ as the number of approximately independent loci $L\to\infty$, with
$\Sigma_0\succ 0$ the per-locus feature covariance under $m=0$. Then the per-locus Fisher
information for the magnitude $m$ at $m=0$ is
\[
\mathcal{I}_0=\Delta_\beta^\top\,\Sigma_0^{-1}\,\Delta_\beta,
\]
and any regular estimator $\widehat m$ obeys the Cram\'er--Rao bound
$\operatorname{Var}(\widehat m)\ge 1/(L\,\mathcal{I}_0)$. Consequently the best achievable
relative error satisfies
\[
\frac{\sqrt{\operatorname{Var}(\widehat m)}}{m}\;\ge\;\frac{1}{m\,\sqrt{L\,\mathcal{I}_0}}
\;\xrightarrow[m\to 0]{}\;\infty,
\]
so at fixed genome size the relative error of magnitude estimation necessarily diverges as the
migration rate vanishes.
\end{thm}

\begin{proof}
By Theorem~\ref{thm:expansion}, $\mathbb{E}_m\Phi=\mathbb{E}_0\Phi+m\,\Delta_\beta+O(m^2)$, so
the limiting Gaussian family $\mathcal{N}\!\big(\mu(m),\,\Sigma_0/L\big)$ has mean
differentiable in $m$ with $\mu'(0)=\Delta_\beta$. For a Gaussian location family with known
covariance the Fisher information for the scalar $m$ is
$\mu'(0)^\top(\Sigma_0/L)^{-1}\mu'(0)=L\,\Delta_\beta^\top\Sigma_0^{-1}\Delta_\beta$. The
Cram\'er--Rao inequality (equivalently the H\'ajek--Le~Cam local asymptotic minimax bound for
this LAN family) gives $\operatorname{Var}(\widehat m)\ge 1/(L\,\mathcal{I}_0)$ for regular
estimators. Dividing the resulting standard-deviation floor $1/\sqrt{L\,\mathcal{I}_0}$ by $m$
and letting $m\to 0$ yields the stated divergence.
\end{proof}

\begin{cor}[Appreciable-rate threshold and the gate]\label{cor:appreciable-threshold}
The relative-error floor $1/(m\sqrt{L\,\mathcal I_0})$ of Theorem~\ref{thm:fisher} crosses one at
\[
   m_\star=\frac{1}{\sqrt{L\,\mathcal I_0}}=\frac{1}{\sqrt{L}\,\lVert\Sigma_0^{-1/2}\Delta_\beta\rVert}.
\]
For $m\ll m_\star$ the standard deviation of every regular estimator exceeds the estimand, so no
scale-accurate magnitude estimate exists; for $m\gg m_\star$ the floor drops below one and magnitude
becomes estimable to $O(1)$ relative error. This is the information-theoretic origin of the
appreciable regime observed empirically: the two-stage procedure that first tests whether
$m\gtrsim m_\star$ and only then estimates or orients is not a heuristic but the decomposition
dictated by the Fisher geometry, which is why direction and magnitude are recoverable exactly in the
appreciable band and near-chance below it. The threshold scales as $m_\star\propto L^{-1/2}$, so a
larger genome lowers it, but only at the slow $L^{-1/2}$ rate; the ordering of replicates by rate,
which depends on the sign of $m-m_\star$ rather than on resolving $m$ itself, survives well below
$m_\star$ and is the ordinal signal the regressor still recovers.
\end{cor}

\begin{prop}[Two thresholds: detection and magnitude versus orientation]\label{prop:two-thresholds}
In the local asymptotic normality regime of Proposition~\ref{prop:linear-optimal}, write
$W=\Sigma_0^{-1/2}$ for the whitening map. Two distinct information scales govern the pipeline.
\emph{Detection and magnitude} depend on the individual first-order shift through
$\lVert W\Delta_\beta\rVert=\sqrt{\mathcal I_0}$: a rate-$m$ replicate sits at Mahalanobis distance
$m\sqrt{L}\,\lVert W\Delta_\beta\rVert$ from the no-migration control, so the detection error is
$Q(\tfrac12 m\sqrt{L}\,\lVert W\Delta_\beta\rVert)$ and the threshold is
$m_\star=1/(\sqrt{L}\,\lVert W\Delta_\beta\rVert)$ (Corollary~\ref{cor:appreciable-threshold}).
\emph{Orientation} depends instead on the pairwise \emph{differences}
$\lVert W(\Delta_\beta-\Delta_{\beta'})\rVert$: by Proposition~\ref{prop:linear-optimal} the
three-way direction accuracy is framed by
\begin{align*}
1-\!\!\sum_{\beta'\neq\beta}\!Q\!\big(\tfrac12 m\sqrt{L}\,
  \lVert W(\Delta_\beta-\Delta_{\beta'})\rVert\big)
&\ \le\ A_\beta(m) \\
&\ \le\ 1-\max_{\beta'\neq\beta}Q\!\big(\tfrac12 m\sqrt{L}\,
  \lVert W(\Delta_\beta-\Delta_{\beta'})\rVert\big),
\end{align*}
so it is controlled by the \emph{smallest} pairwise separation, with orientation threshold
\[
m_\star^{\mathrm{dir}}
=\frac{c}{\sqrt{L}\,\min_{\beta\neq\beta'}
  \lVert W(\Delta_\beta-\Delta_{\beta'})\rVert}.
\]
Both scale as $L^{-1/2}$ but are set by different Fisher functionals---the individual shift for
detection, the pairwise differences for orientation---which is why the gate can flag appreciable
flow at rates where the direction is not yet resolvable.
\end{prop}

\begin{proof}
The detection statement is the overlap of $\mathcal N(\mu_0,\Sigma_0/L)$ with
$\mathcal N(\mu_0+m\Delta_\beta,\Sigma_0/L)$, whose Mahalanobis separation is
$m\sqrt{L}\,\lVert W\Delta_\beta\rVert$; equating the resulting standard-deviation floor to $m$
recovers Corollary~\ref{cor:appreciable-threshold}. For orientation, the class-$\beta$ region of the
Bayes rule is the intersection of the pairwise half-spaces of
Proposition~\ref{prop:linear-optimal}; a union bound over the two competitors gives the lower bound
and retaining the single nearest competitor the upper bound. Both are monotone in
$\min_{\beta'}\lVert W(\Delta_\beta-\Delta_{\beta'})\rVert$, which therefore sets the threshold.
\end{proof}

\begin{thm}[All ordered directions: logarithmic class-count cost and permutation equivariance]
\label{thm:many-directions}
Let $s\ge2$ sampled populations be represented by the population-labelled rarefaction map
$\Phi_s$ containing every $\alpha_i$, every $\pi_i$, and every pair-private
$\widehat\pi_{ij}$, with the same across-locus moments and depth summaries used above. There are
$K=s(s-1)$ ordered donor--recipient edges $\mathcal B_s=\{i\!\to\!j:i\ne j\}$. Assume the local
Gaussian experiment of Proposition~\ref{prop:linear-optimal}, now for
$\beta\in\mathcal B_s$, with equal class priors, common positive-definite covariance
$\Sigma_0/L$ on the identifiable feature span, and
$\mu_\beta(m)=\mu_0+m\Delta_\beta+O(m^2)$. Put
\[
 d_{\beta\beta'}=m\sqrt L\,
 \big\|\Sigma_0^{-1/2}(\Delta_\beta-\Delta_{\beta'})\big\|,
 \qquad
 \gamma_s=\min_{\beta\ne\beta'}
 \big\|\Sigma_0^{-1/2}(\Delta_\beta-\Delta_{\beta'})\big\|.
\]
Then, to first order in $m$:
\begin{enumerate}
\item the $K$-class Bayes rule is the linear discriminant
\[
 \widehat\beta(\phi)=\arg\max_{\beta\in\mathcal B_s}
 \left\{\mu_\beta(m)^\top\Sigma_0^{-1}\phi
 -\tfrac12\mu_\beta(m)^\top\Sigma_0^{-1}\mu_\beta(m)\right\};
\]
\item for a true edge $\beta$, its conditional accuracy $A_\beta$ satisfies
\[
 \max\!\left\{0,1-\!\sum_{\beta'\ne\beta}Q(d_{\beta\beta'}/2)\right\}
 \le A_\beta
 \le 1-\max_{\beta'\ne\beta}Q(d_{\beta\beta'}/2);
\]
consequently, for $0<\varepsilon<(K-1)/2$, the sufficient condition
\begin{equation}
 m\ \ge\ \frac{2}{\sqrt L\,\gamma_s}\,
 Q^{-1}\!\left(\frac{\varepsilon}{K-1}\right)
 \label{eq:many-direction-threshold}
\end{equation}
gives $A_\beta\ge1-\varepsilon$ for every edge. If $\gamma_s$ stays bounded away from zero,
the extra standardized separation needed for more labels is only
$Q^{-1}(\varepsilon/(K-1))=\sqrt{2\log((K-1)/\varepsilon)}\,(1+o(1))$, hence
$O(\sqrt{\log K})$, not $O(K)$;
\item for each population relabelling $\rho\in S_s$ there is a coordinate-permutation matrix
$R_\rho$ such that $\Phi_s(\rho\cdot G)=R_\rho\Phi_s(G)$. If the base demography and edge prior
are invariant under $S_s$, then $\Delta_{\rho\cdot\beta}=R_\rho\Delta_\beta$ and
$\Sigma_0=R_\rho\Sigma_0R_\rho^\top$. Moreover any probabilistic classifier $q$ can be
symmetrized to an equivariant classifier
\[
 \overline q(\beta\mid\phi)=\frac1{|S_s|}\sum_{\rho\in S_s}
 q(\rho^{-1}\!\cdot\beta\mid R_{\rho^{-1}}\phi)
\]
whose population cross-entropy risk is no larger than that of $q$.
\end{enumerate}
\end{thm}

\begin{proof}
For common-covariance Gaussian classes the log likelihood is, up to a term independent of
$\beta$, the affine score in part~1. Under class $\beta$, the score difference between a
competitor $\beta'$ and $\beta$ is Gaussian with mean $-d_{\beta\beta'}^2/2$ and variance
$d_{\beta\beta'}^2$. Thus the probability that this one competitor wins is exactly
$Q(d_{\beta\beta'}/2)$. Multiclass error is the union of these $K-1$ competitor events: its
probability is at least their largest probability and at most their sum. Taking complements
proves part~2. Since every distance is at least $m\sqrt L\,\gamma_s$, the union bound is at most
$(K-1)Q(m\sqrt L\,\gamma_s/2)$; solving for an error at most $\varepsilon$ gives
\eqref{eq:many-direction-threshold}, and the
normal-tail inverse has the stated small-argument asymptotic.

The rarefaction formulas are unchanged by renaming populations: relabelling only permutes the
$\alpha_i$, $\pi_i$, and $\widehat\pi_{ij}$ coordinates, and the moment and depth operators commute
with that permutation. This proves the first assertion of part~3; invariance of the base law then
gives the response and covariance identities by differentiation at $m=0$. The displayed group
average is equivariant by a change of variables in $S_s$. Finally, $z\mapsto-\log z$ is convex,
so Jensen's inequality bounds the loss of $\overline q$ by the average loss of its $|S_s|$
relabelled copies. Invariance of the joint law of $(\Phi_s,\beta)$ makes every term in that average
equal to the risk of $q$.
\end{proof}

\begin{rem}[Numerical validation and the detectability ordering]\label{rem:threshold-numerics}
Estimating $\mu_\beta$ and $\Sigma_0$ directly from the replicate features, the Gaussian rule of
Proposition~\ref{prop:linear-optimal} reproduces the empirical per-rate direction accuracy across
three orders of magnitude in the migration rate---$0.38/0.40/0.69/1.00$ predicted against
$0.34/0.34/0.74/1.00$ measured (linear discriminant, replicate-aggregated) at
$5\times10^{-7}/2.5\times10^{-6}/5\times10^{-5}/2.5\times10^{-4}$---so the threshold formula is the
operating curve, not an analogy. The whitened shifts order the three scenarios by detectability,
$\lVert W\Delta_C\rVert>\lVert W\Delta_B\rVert>\lVert W\Delta_A\rVert$
($1.73:1.30:0.62$ at $5\times10^{-5}$): introgression from the divergent outgroup into a sister
($C=P_3\!\to\!P_2$) leaves the largest rarefaction footprint, and gene flow between the two similar
sisters ($A=P_1\!\to\!P_2$) the smallest. The binding constraint on orientation is therefore the
sister-to-sister scenario $A$, not the reversed-flow pair $B$--$C$: the residual confusion is largest
for $A$ and smallest for $B$--$C$ (pair confusions $307\,{:}\,217\,{:}\,160$ over all positive rates,
Figure~\ref{fig:class-confusion-detailed}), so the population-labelled channels resolve the
donor--recipient orientation of the reversed pair \emph{best}---precisely the call $D$ cannot
make---while the residual difficulty is the weak footprint of gene flow between the two most similar
populations.
\end{rem}

\begin{thm}[Row-level split optimism under replicate random effects]\label{thm:leakage}
Suppose replicate \(i\) contributes rows
\[
X_{i,g}=\mu(Y_i,m_i,g)+Z_i+\varepsilon_{i,g},\qquad g=2,\ldots,199,
\]
where \(Z_i\) is a replicate-level random effect shared by all rarefaction depths from the
same simulation. Under an independent row-level train/test split with test probability
p, the probability that replicate \(i\) appears in both training and test is
\[
1-(1-p)^{198}-p^{198}.
\]
For \(p=0.2\) this probability is essentially one. Consequently row-level test risk
estimates within-replicate interpolation risk rather than new-replicate generalization
risk. If the learner can estimate \(Z_i\), the row-level risk is optimistically biased by
a nonnegative term increasing with \(\mathrm{Var}(Z_i)\).
\end{thm}

\begin{proof}
The overlap probability is the complement of the event that all 198 rows fall in training
or all 198 rows fall in test. This gives the displayed expression. In a squared-error
model, a learner that sees training rows from replicate \(i\) can estimate the shared
effect \(Z_i\) and use it when predicting held-out rows from the same replicate. For a new
replicate, no such estimate is available. The conditional test error under row splitting
therefore removes part of the replicate-level variance that remains present in
new-replicate prediction. The same argument applies to classification through any margin
model in which estimating \(Z_i\) moves examples away from the decision boundary. Hence
row-level splitting is optimistic whenever replicate effects carry predictive signal.
\end{proof}

\begin{thm}[Concentration of empirical rarefaction features]
Assume \(L\) independent loci and let \(\widehat{\Phi}\in\mathbb{R}^d\) be the empirical
rarefaction feature vector, with each coordinate bounded in an interval of length \(B\).
Then
\[
\Pr\left(\|\widehat{\Phi}-\mathbb{E}\Phi\|_\infty>\varepsilon\right)
\le 2d\exp\left(-\frac{2L\varepsilon^2}{B^2}\right).
\]
\end{thm}

\begin{proof}
Apply Hoeffding's inequality to each bounded coordinate of the empirical average over
loci. A union bound over the \(d\) coordinates gives the displayed inequality. For linked
loci the same proof applies after replacing \(L\) by an effective number of approximately
independent recombination blocks.
\end{proof}

\begin{rem}
This concentration is the finite-$L$ basis for the local asymptotic normality invoked in
Theorem~\ref{thm:fisher} and Proposition~\ref{prop:linear-optimal}: the empirical feature vector
sits within $O(L^{-1/2})$ of its expectation, so at fixed genome size the recoverable signal is the
first-order mean shift $m\Delta_\beta$ measured against an $L^{-1/2}$ noise floor. The ratio of the
two is what the Fisher threshold $m_\star$ of Corollary~\ref{cor:appreciable-threshold} balances,
and it is the same quantity that sets both the direction Bayes error and the magnitude relative
error---one information-geometric scale governing the whole pipeline.
\end{rem}

\section{Discussion}
\label{sec:discussion}

The contribution of this paper is a compact, interpretable, outgroup-free route to the
\emph{direction} of introgression. Under strict replicate-aware evaluation---no simulation
replicate shared between training and test---DNNaic orients appreciable gene flow at $96.6\%$ at
the largest fixed rate and $99.7\%$ across the appreciable band, with a detect-then-orient gate that
flags appreciable migration at ROC-AUC $0.99$ and a well-calibrated score; accuracy falls to the
chance rate as migration becomes negligible, as the identifiability theory predicts. On the same
simulations a precomputed Patterson's $D$ labels direction below the majority-class baseline and
barely detects presence---not incidental, but what Theorem~\ref{thm:symmetric} forces, since $D$
collapses the population labels that carry orientation. The signal lives in the asymmetric,
population-labeled private and pair-private channels (Table~\ref{tab:channels})---precisely the
components that symmetric, label-collapsing $D$-statistics discard (Theorem~\ref{thm:symmetric}).

The all-edges experiment shows both how the method scales and what kind of structure pays. With
four exchangeable populations, the exact 12-way task reaches $62.9\%$ versus $8.3\%$ chance when
the $S_4$ relabelling symmetry is enforced inside training folds, and the known-pair reversal task
reaches $78.8\%$ versus $50\%$. Removing population labels or permuting replicate labels returns
accuracy to chance. Neither nonlinear kernels nor neural and random-matrix-whitened learned
representations improve on the linear rule. This is the empirical counterpart of
Theorem~\ref{thm:many-directions}: the useful leverage is equivariance of the labelled feature map,
not unconstrained capacity. Just as importantly, the frozen nuisance audit shows that class-count
scaling does not imply demographic portability; weaker signal, unequal sizes, and topology shifts
have to be represented in training.

A distinguishing feature of this work is that each theoretical claim is paired with a numerical
check on the same arrays. The symmetric-summary obstruction (Theorem~\ref{thm:symmetric}) is borne
out directly: a classifier restricted to the symmetric richness statistics reaches only $80.2\%$ in
the appreciable band, and it is the asymmetric private and pair-private channels---the ones the
theorem singles out---that lift it to $99.8\%$, while a precomputed Patterson's $D$, which discards
those labels, cannot orient at all. The linear-optimality prediction
(Proposition~\ref{prop:linear-optimal}) is confirmed by a logistic model matching the network
($78.0\%$ versus $75.6\%$ overall). The replicate-effect leakage theorem (Theorem~\ref{thm:leakage})
is confirmed by the row-versus-group gap: the same features score $\approx99.6\%$ under a leaky
row-level split (Section~\ref{subsec:benchmarks}) but $\approx75\%$ under the correct replicate
split. And the Fisher-information limit (Theorem~\ref{thm:fisher}) is matched by the divergence of
the magnitude relative error toward small rates (Section~\ref{sec:regression_model}). This
theory--numerics correspondence is what turns an empirical neural-network fit into a
population-genetic inference method, and it positions DNNaic in the summary-statistic learning
tradition
\citep{sheehan2016deep,durvasula2019archie} rather than the genotype-matrix CNN tradition
\citep{gower2021genomatnn,ray2024introunet,flagel2019unreasonable,chan2018likelihood}.  ArchIE
\citep{durvasula2019archie} is the closest analogue: it feeds a fixed set of summary statistics
to a simple learner, and DNNaic's higher-moment rarefaction features extend that idea to the direction
of gene flow.

The direction result is established under ground-truth replicate identity
(Section~\ref{subsec:true-replicate}) and reproduced independently across regenerations and two
model architectures, confirming high-accuracy recovery wherever migration is appreciable and a
principled decline to chance where it is not. The magnitude component, shown here to be
information-limited rather than mis-modeled, is reported per rate bin and paired with a power
analysis in the number of independent loci. The \textit{Heliconius} audit separates two claims that
the initial positive panel had conflated: Patterson's $D$ and a raw pair-private contrast track the
selected geographic signal relative to an allopatric control, but the frozen direction head calls
the same class on both at high score (Section~\ref{subsec:heliconius}). This failure does not erase
the rarefaction asymmetry; it shows that a simulation-trained decision boundary is not a calibrated
natural-panel orientation rule without matched training and controls.

\section{Limitations and scope}
\label{sec:limitations}
We delimit the scope of the present study.
\begin{itemize}
\item \textbf{Direction, not magnitude.} The simulation-supported deliverable is the direction of
appreciable gene flow within the evaluated generative families. Migration-rate \emph{magnitude} is information-limited at small rates: positive-rate
relative error is dominated by the smallest rates, which Theorem~\ref{thm:fisher} shows to be an
intrinsic Fisher-information limit at fixed genome size---consonant with established limits on
SFS-based inference \citep{terhorst2015fundamental}. Magnitude should therefore be reported with
per-rate breakdowns and treated as reliable only above a rate threshold that shrinks as the number
of independent loci grows.
\item \textbf{Simulation-trained.} Quantitative results are on coalescent simulations under
hierarchical Wright--Fisher demographies. A scoped cross-demography probe
(Section~\ref{subsec:true-replicate}) shows the frozen direction model transfers across alternative
histories---perfectly to a deeper, larger-$N_e$ demography and at $73\%$ (versus $33\%$ chance) to a
recent, small-$N_e$ one---so the orientation signal is not merely memorized from the training history,
though it attenuates when coalescence is fast and divergence shallow. Broader nuisance tests are more
sobering, and the \textit{Heliconius} allopatric control is decisive for interpretation: despite an
approximately null $D$ and a reversed raw sharing ratio, the frozen head still calls $C$ at score
$0.9918$. The natural-data head is therefore uncalibrated, not externally validated. A model matched
to the target taxon and independent labelled controls are prerequisites, not optional refinements.
The complementary limit is visible on the Neanderthal system (Section~\ref{subsec:neanderthal}): there the frozen model
detects the introgression through Patterson's $D$ but \emph{abstains} from orienting it, because only
a few high-coverage archaic genomes exist (rarefaction depth $g\le4$) and the archaic lineage lies
outside the training distribution. That abstention is a useful case, not a general guarantee: the
\textit{Heliconius} control shows that a different off-distribution panel can receive a confident
direction. Reliability therefore has to be established in the target system rather than inferred
from the simulation calibration curve.
\item \textbf{All-edges extension is a controlled benchmark.} The 12-direction result uses an
exchangeable four-tip star, $20$ independent baseline replicates per edge, and a strong integrated
exposure $mT=0.25$; its Wilson interval and three split-seed results are reported rather than a
single point estimate. It establishes that all ordered edges can be learned above chance and that
permutation symmetry is valuable, not that accuracy remains $62.9\%$ for arbitrary trees, sample
sizes, or migration histories. The four one-factor stress tests contain only eight replicates per
edge and are diagnostic. A definitive multi-population model should train over a broad prior on
topology, $N_e$, timing, pulse shape, recombination, missingness, and linked selection, then be
tested on a preregistered independent simulation bank and target-system controls.
\item \textbf{Theory regime.} The perturbation and identifiability results
(Theorems~\ref{thm:expansion}--\ref{thm:fisher}) are first-order in the migration magnitude $m$ and
assume a finite-state truncation of the ancestral process; the direction-separation margin $\gamma$
in Corollary~\ref{cor:identifiability} is established generically, and computing it for a specific
demography by finite differences in \texttt{msprime} is a natural next step.
\end{itemize}

\subsection*{Data and code availability}
\textsc{PADZE}, the software used to compute every rarefaction feature in this study, is
open-source at \url{https://github.com/Andres42611/PADZE}~\citep{padze2026}. The complete
simulation dataset (the $3{,}200$-replicate rarefaction feature arrays with ground-truth
replicate, direction, and rate labels) and the leakage-controlled evaluation code are
archived at a public Zenodo record~\citep{delcastillo2026data}.

\subsection*{Funding}
The research presented in this paper was supported by the European Research Council (ERC) under
the European Union's Horizon 2022 research and innovation programme (grant agreement No.\ 101041711),
by the Simons Foundation, by Heights Labs, by the Israel Science Foundation (grant number 2258/19),
and by the Israel Science Foundation (ISF Grant 4101/25).

\subsection*{Declaration on generative AI}
During the preparation of this work the authors used generative AI tools to accelerate drafting and
revision. All mathematical claims, proofs, and bibliographic information were subsequently reviewed
and edited by the authors, who take full responsibility for the content of the manuscript.

\clearpage
\bibliographystyle{abbrvnat} % or plainnat / unsrtnat
\bibliography{refs}          % refs.bib file name (no .bib extension)

\end{document}
